\chapter{推理优化：让模型跑得更快}
\label{ch:inference}
%% ============================================================

训练完成后，模型需要服务用户。
在推理阶段，挑战从「如何训练更好的模型」
转变为「如何让好模型跑得更快、更便宜」。

\section{推理的根本瓶颈}

在讨论具体的优化技术之前，
我们需要先理解 LLM 推理的\textbf{根本瓶颈}，
只有抓住瓶颈，才能理解为什么每种优化手段有效。

LLM 推理分为两个截然不同的阶段：

\textbf{Prefill 阶段}（首 token 生成前）：
模型接收用户的完整输入（prompt），
\textbf{并行}计算所有输入 token 的 KV 值。
这个阶段是\textbf{计算瓶颈}的：
因为有大量 token 需要同时处理，
GPU 的计算单元（CUDA core / Tensor Core）被充分利用。
Prefill 的耗时决定了 \textbf{TTFT}（Time To First Token，首 token 延迟）。

\textbf{Decode 阶段}（逐 token 生成）：
模型每次只生成\textbf{一个} token，
但每生成一个 token，都需要读取\textbf{整个模型的权重}。
这个阶段是\textbf{内存带宽瓶颈}的，
GPU 的计算能力远远没有被用满，
大部分时间花在等待从 HBM（高带宽显存）读取权重。
Decode 的速度决定了 \textbf{ITL}（Inter-Token Latency，token 间延迟）。

\begin{whybox}[为什么 Decode 是内存带宽瓶颈？]
以一个 70B 参数的模型（BF16 精度）为例。
每生成一个 token，需要读取全部权重：
$70 \times 10^9 \times 2\,\text{bytes} = 140\,\text{GB}$。
但一个 token 的计算量只有约 $70 \times 10^9$ 次乘加运算（multiply-add）。

H100 GPU 的 HBM 带宽是 3.35\,TB/s，
读取 140\,GB 需要 $140 / 3350 \approx 42\,\text{ms}$。
H100 的 BF16 算力是 989\,TFLOPS，
$70 \times 10^9$ 次运算只需要 $0.07\,\text{ms}$。

\textbf{计算比内存读取快 600 倍}，
GPU 99.8\% 的时间在等待数据从显存搬运到计算单元。
这就是为什么 LLM decode 阶段是内存带宽瓶颈的。
\end{whybox}

理解了这个瓶颈，一个关键的杠杆就浮现出来：\textbf{batch size}。
如果同时处理 $B$ 个请求，
每次读取权重后可以为 $B$ 个 token 做计算。
读取成本不变，但有效计算量增加了 $B$ 倍。

存在一个\textbf{临界 batch size}：
\begin{equation}
  B_{\text{critical}} = \frac{\text{计算能力 (FLOPS)}}{\text{内存带宽 (bytes/s)} \times \text{算术强度}}
\end{equation}
当 $B < B_{\text{critical}}$ 时，推理受内存带宽限制；
当 $B > B_{\text{critical}}$ 时，推理变为计算瓶颈。
对于 H100，临界 batch size 大约在 150--300 之间
（取决于模型架构和序列长度）。

在生产环境中，还需要在三个指标之间做权衡：
\begin{itemize}[noitemsep]
  \item \textbf{TTFT}（首 token 延迟）：用户等待第一个字出现的时间。
        对交互体验影响最大。取决于 prefill 速度。
  \item \textbf{ITL}（token 间延迟）：两个输出 token 之间的间隔。
        决定了「打字速度」感受。取决于 decode 速度。
  \item \textbf{Throughput}（吞吐量）：每秒处理的总 token 数。
        决定了服务成本。可以通过增大 batch size 提升，
        但会牺牲 TTFT 和 ITL。
\end{itemize}

这三个指标往往相互矛盾：
增大 batch size 提升吞吐量，但增加了每个请求的延迟。
推理框架的核心任务就是在这三者之间找到最优平衡。

理解这些指标的\textbf{量级}有助于建立正确的直觉。
以 Llama 3.1 70B（FP16）在单个 H100 上的推理为例：
\begin{itemize}[noitemsep]
  \item \textbf{TTFT}（batch=1，prompt 1024 tokens）：
        $\sim$120\,ms。主要时间花在 prefill 计算上。
  \item \textbf{ITL}（batch=1）：
        $\sim$42\,ms（$\approx 140\,\text{GB} / 3.35\,\text{TB/s}$）。
        几乎完全由读取权重的时间决定。
  \item \textbf{ITL}（batch=64）：
        $\sim$45\,ms。只比 batch=1 慢 7\%，
        因为额外的 63 个请求的计算几乎不增加延迟
        （计算量仍在带宽瓶颈之下）。
  \item \textbf{Throughput}（batch=64）：
        $\sim$64 / 0.045 $\approx$ 1,422 tok/s。
        而 batch=1 只有 $\sim$24 tok/s。
        \textbf{batch size 增加 64$\times$，
        吞吐量增加 $\sim$60$\times$}，
        这就是为什么高效的 batching 是推理框架的核心。
\end{itemize}

量化将权重从 FP16 减半到 FP8 后，
上述所有数字几乎\textbf{翻倍}（内存读取减半 → 速度翻倍）。
INT4 量化则接近 $4\times$ 加速，
这解释了为什么量化是推理优化中
投入产出比最高的技术手段。

\begin{keybox}[推理成本的“\$1/M tokens”基准（2026 年初）]
在 H200 云实例（$\sim$\$3.5/hour）上部署不同模型的参考成本：
\begin{itemize}[noitemsep]
  \item \textbf{Llama 3.1 8B}（FP8）：$\sim$\$0.03/M tokens。
        $\sim$12,000 tok/s 吞吐。日常任务的经济之选。
  \item \textbf{Llama 3.1 70B}（FP8）：$\sim$\$0.25/M tokens。
        需要 2$\times$ H200（张量并行），$\sim$3,000 tok/s。
  \item \textbf{DeepSeek-V3/R1}（FP8，MoE）：$\sim$\$0.15/M tokens。
        需要 8$\times$ H200（专家并行），$\sim$2,200 tok/s/GPU。
        得益于 MoE 的活跃参数效率，
        单位成本反而低于同等效果的 70B 稠密模型。
  \item \textbf{Llama 3.1 8B}（Q4\_K\_M，MacBook M4）：\$0（本地运行）。
        $\sim$40 tok/s，适合隐私敏感的个人使用。
\end{itemize}
\end{keybox}


\section{量化：精度换速度}

量化的核心思想是：
模型推理时的大部分计算不需要训练时那么高的数值精度。
将权重和/或激活从高精度（FP16/BF16）
转换为低精度（INT8、INT4 甚至 FP4），
可以显著减少内存占用和计算时间。

\subsection{量化基础}

量化的本质是将连续（或高精度离散）的数值
映射到有限数量的离散级别（discrete levels）。
对称线性量化公式：
\begin{equation}
  x_q = \mathrm{round}\!\left(\frac{x}{s}\right), \quad
  s = \frac{\max(|x|)}{2^{b-1} - 1}
\end{equation}
其中 $s$ 是缩放因子，$b$ 是目标位宽，
$x_q$ 是量化后的整数值。
解量化时，$\hat{x} = x_q \cdot s$。

非对称量化增加了一个零点（zero-point）参数 $z$：
\begin{equation}
  x_q = \mathrm{round}\!\left(\frac{x - x_{\min}}{s}\right), \quad
  s = \frac{x_{\max} - x_{\min}}{2^b - 1}, \quad
  z = \mathrm{round}\!\left(\frac{-x_{\min}}{s}\right)
\end{equation}
非对称量化能更好地处理分布不对称的权重，
但增加了计算开销（解量化需要额外的加法）。

常用的数值格式各有优劣：
\begin{itemize}[noitemsep]
  \item \textbf{INT8}：8-bit 整数。256 个级别，覆盖 $[-128, 127]$。
        精度损失几乎为零，是最安全的量化选择。
  \item \textbf{FP8（E4M3/E5M2）}：8-bit 浮点。
        E4M3 有更高的精度（3-bit 尾数），E5M2 有更大的动态范围（5-bit 指数）。
        H100 原生支持 FP8 运算。
  \item \textbf{INT4}：4-bit 整数。仅 16 个级别，
        需要精巧的量化策略来控制精度损失。
  \item \textbf{NF4}（Normal Float 4）：QLoRA~\cite{dettmers2023qlora} 提出的 4-bit 格式。
        16 个量化级别不是均匀分布的，
        而是按照标准正态分布的分位数分布：
        因为预训练后的权重近似正态分布，
        这种非均匀量化能最大化信息保留。
  \item \textbf{NVFP4}：NVIDIA Blackwell 架构原生支持的 4-bit 浮点格式。
        采用 1-bit 符号 + 2-bit 指数 + 1-bit 尾数的结构，
        相比 INT4 保留了更好的动态范围。
        B200 GPU 的 NVFP4 Tensor Core 算力
        达到 FP8 的 $2\times$，
        使得 4-bit 推理在数据中心场景中首次变得实用。
  \item \textbf{MXFP8}（Microscaling FP8）：
        AWS Trainium3 采用的新一代 8-bit 浮点标准。
        与传统 FP8 的关键区别在于
        MXFP8 引入了\textbf{微缩放}（microscaling）机制，
        在细粒度的元素组（通常 32 个元素）级别
        共享一个缩放因子，
        相比 per-tensor 缩放提供了更好的精度保持。
        Trainium3 原生支持 MXFP8 用于训练和推理，
        标志着 FP8 变体正在从 NVIDIA 专属格式
        走向跨平台标准化。
  \item \textbf{ExpertsInt8}（AI21 Jamba）：
        专为 MoE 架构设计的量化方案。
        核心洞察是 MoE 模型中\textbf{不同专家}的权重分布差异巨大，
        传统的统一量化策略无法很好地适应。
        ExpertsInt8 为每个专家独立计算量化参数，
        使得 Jamba 的 94B 活跃参数 MoE 模型
        能够在单台机器上运行，
        大幅降低了 MoE 推理的硬件门槛。
\end{itemize}

\subsection{Post-Training Quantization（PTQ）}

PTQ 在训练完成后直接对权重进行量化，无需重新训练。
最简单的方法是 Round-to-Nearest（RTN），
直接将每个权重四舍五入到最近的量化级别。
RTN 对于 INT8 通常效果不错，
但在 INT4 及以下位宽时精度损失明显。
原因是量化误差在层间\textbf{累积}，
一层的量化误差会影响下一层的输入分布，
导致误差逐层放大。

\textbf{GPTQ}~\cite{frantar2022gptq} 使用二阶信息（Hessian 矩阵）来最小化量化误差。
核心思想是：量化一个权重时，
利用 Hessian 信息调整同一行中\textbf{尚未量化}的权重，
来补偿量化引入的误差。
GPTQ 逐列处理权重矩阵：
量化第 $i$ 列后，根据 Hessian 的逆矩阵
调整第 $i+1$ 到第 $n$ 列的未量化权重。
这使得整行的\textbf{总输出}
尽可能接近未量化时的输出。
GPTQ 通常只需要一小批校准数据（128 条样本）
和几分钟的计算时间，就能完成整个模型的量化。

具体而言，GPTQ 源自 OBQ（Optimal Brain Quantization）算法。
对于权重矩阵 $W$ 的每一行 $w$，
目标是找到量化后的 $\hat{w}$ 使得输出误差最小：
\begin{equation}
  \min_{\hat{w}} \| W x - \hat{W} x \|^2
  = \min_{\hat{w}} (w - \hat{w})^T H (w - \hat{w})
\end{equation}
其中 $H = X X^T$ 是校准数据上的 Hessian 矩阵。
当量化第 $i$ 列时，剩余列的最优调整量为：
\begin{equation}
  \delta_F = -\frac{w_i - \mathrm{quant}(w_i)}{H_{ii}^{-1}} \cdot (H^{-1})_{:,i}
\end{equation}
这个公式的含义是：根据 Hessian 的逆（衡量每个权重的敏感度），
将第 $i$ 列的量化误差\textbf{分摊}到后续未量化的列上。
敏感度低的权重承担更多补偿，敏感度高的权重保持更精确。

GPTQ 的关键创新在于发现\textbf{列的处理顺序}
对最终精度的影响远小于预期，
这使得可以固定一个处理顺序（而不是贪心搜索最优顺序），
将时间复杂度从 $O(d^3 \cdot n)$ 降低到 $O(d \cdot n^2)$，
使得量化一个 175B 的模型只需要约 4 GPU-hours。

\textbf{AWQ}（Activation-Aware Weight Quantization）~\cite{lin2024awq}
发现了一个关键洞察：
\textbf{并非所有权重同等重要}。
少数权重（约 1\%）对模型输出的影响远超其他权重，
这些权重对应的激活值特别大。
AWQ 通过激活值的大小来识别这些「显著权重」（salient weights），
对它们施加更大的缩放因子以保护精度，
同时让不重要的权重承受更大的量化误差。

AWQ 的核心操作是引入一个 per-channel 的缩放因子 $s$：
\begin{equation}
  Q(w \cdot s) \cdot \frac{x}{s} \approx w \cdot x
\end{equation}
对于显著通道（激活值大的通道），$s > 1$，
使得权重被放大后量化精度更高；
对于不重要的通道，$s \leq 1$，允许更大的量化误差。
缩放因子 $s$ 通过在校准数据上搜索最优值来确定。

AWQ 在 INT4 量化下通常优于 GPTQ，
因为它把有限的精度预算花在了刀刃上。
在实践中，AWQ 搭配 Marlin 高性能 kernel
可以在 H200 上达到 741 tok/s（Qwen2.5-32B），
超过 FP16 的 461 tok/s。

\textbf{SmoothQuant}~\cite{xiao2023smoothquant} 解决的是\textbf{激活值量化}的难题。
权重的分布通常比较平滑，容易量化；
但激活值中存在大量\textbf{离群值}（outliers），
个别通道的值可能比其他通道大 100 倍。
这些离群值让量化极为困难：
为了覆盖最大值，缩放因子必须设得很大，
导致其他正常值的量化精度严重降低。

SmoothQuant 的解法很优雅：
在矩阵乘法 $Y = XW$ 中，引入一个对角矩阵 $\mathrm{diag}(s)$：
\begin{equation}
  Y = XW = (X \cdot \mathrm{diag}(s)^{-1}) \cdot (\mathrm{diag}(s) \cdot W)
  = \hat{X} \hat{W}
\end{equation}
选择合适的 $s$，
可以将激活中的离群值「转移」到权重上，
$\hat{X}$ 变得平滑容易量化，
$\hat{W}$ 虽然变大了一些，但权重本身分布均匀，仍然可以量化。
数学上这是等价变换，不改变计算结果。

\textbf{SqueezeLLM}~\cite{kim2024squeezellm}
从一个不同的角度切入量化问题：
\textbf{灵敏度感知}（sensitivity-aware）的非均匀量化。
SqueezeLLM 的核心观察是，
权重矩阵中存在少数\textbf{灵敏度极高}的离群权重，
这些权重虽然数量稀少（通常不到 0.5\%），
但对输出质量的影响不成比例地大。
均匀量化对所有权重一视同仁，
必然在离群值和普通值之间做出糟糕的折中。

SqueezeLLM 的解决方案分两步：
首先，通过 Fisher 信息矩阵（近似 Hessian 对角线）
识别出灵敏度最高的权重，
将它们从权重矩阵中\textbf{抽离}出来，
以全精度稀疏格式单独存储；
然后，对剩余的权重做\textbf{非均匀量化}，
量化级别不再是等间距分布的，
而是通过 k-means 聚类找到能最小化
加权重构误差的最优量化码本。
这种“分离 + 非均匀”策略在 INT3 量化下
就能达到其他方法 INT4 的精度水平，
将压缩率推向了新的极限。

\textbf{FP8 量化：从推理到训练的统一}
是 2024--2026 年数据中心量化的核心趋势。
H100 是第一款原生支持 FP8 Tensor Core 的 GPU，
它同时支持两种 FP8 变体：
E4M3（4-bit 指数 + 3-bit 尾数，精度高但动态范围小）
和 E5M2（5-bit 指数 + 2-bit 尾数，动态范围大但精度低）。
实践中的最佳策略是\textbf{混合使用}：
前向传播的权重和激活用 E4M3（精度优先），
反向传播的梯度用 E5M2（动态范围优先，
因为梯度值的变化范围更大）。

FP8 训练和 FP8 推理的关键区别在于\textbf{精度要求}：
\begin{itemize}[noitemsep]
  \item \textbf{FP8 推理}（PTQ）：
        直接将训练好的 BF16 权重转换为 FP8，
        通过校准数据确定最优缩放因子。
        精度损失极小（$<$0.5\%），是数据中心的标准做法。
        vLLM 的 \texttt{-{}-dtype fp8} 参数一键启用。
  \item \textbf{FP8 训练}：
        在整个训练过程中使用 FP8 精度。
        需要精心设计的\textbf{损失缩放}（loss scaling）策略
        来防止下溢（FP8 E4M3 的最小非零值约为 $10^{-9}$，
        远大于 BF16 的 $10^{-38}$）。
        DeepSeek-V3 使用 FP8 精度完成了大部分训练，
        证明了 FP8 训练在千亿参数规模上的可行性。
  \item \textbf{FP8 vs INT8 推理}：
        FP8 保留了浮点数的动态范围，
        对激活中的离群值有更好的适应性，
        不需要像 SmoothQuant 那样做激活平滑预处理。
        但 INT8 的硬件支持更广泛（几乎所有现代加速器），
        在不含激活离群值的层中 INT8 和 FP8 精度相当。
\end{itemize}

从 H100 的 FP8 到 Blackwell 的 NVFP4，
再到 Trainium3 的 MXFP8，
硬件厂商正在将低精度计算\textbf{内化为芯片能力}。
这意味着量化从一个“需要精心调参的优化技术”
转变为“硬件原生支持的标准操作”，
2026 年的推理工程师更多的时间花在选择正确的精度格式上，
而不是调整量化超参数。

\subsection{Quantization-Aware Training（QAT）}

QAT 在训练过程中模拟量化效果：
前向传播时使用量化后的权重计算，
反向传播时用直通估计器（Straight-Through Estimator, STE）
近似 round 操作的梯度（round 函数本身梯度为零）。
模型在训练中就学会了适应量化误差，
因此 QAT 的精度通常优于 PTQ，
但代价是需要额外的训练步骤（通常是原始训练的 5--10\%）。

STE 的数学形式很简单：
\begin{equation}
  \frac{\partial \mathcal{L}}{\partial w} \approx
  \frac{\partial \mathcal{L}}{\partial \hat{w}} \cdot 1
  = \frac{\partial \mathcal{L}}{\partial \hat{w}}
\end{equation}
其中 $\hat{w} = \mathrm{round}(w/s) \cdot s$ 是量化后的权重。
$\mathrm{round}$ 的真实梯度为零（几乎处处），
STE 将其替换为恒等函数的梯度（1），
使得梯度可以“穿过”量化操作回传到原始权重。
这个近似虽然粗糙，但在实践中非常有效，
因为在训练的后期，权重的变化量很小，
量化误差的方向大致是随机的，
STE 的偏差在期望意义上趋近于零。

对于 INT8 和 FP8，PTQ 已经足够好，
QAT 的收益不大（通常 $<$0.1\% 的额外精度提升）；
但对于 INT4 及以下位宽，QAT 的优势变得明显。
LLM-QAT~\cite{liu2024llmqat} 表明，
在 INT4 量化下，QAT 相比 GPTQ（PTQ）
可以将 perplexity 增量降低 30--50\%，
特别是在长序列生成任务中差距更大：
因为量化误差在自回归生成中逐 token 累积，
QAT 训练的模型对这种累积误差更具鲁棒性。

QAT 的主要挑战在于\textbf{训练成本}。
对一个 70B 模型做完整的 QAT 通常需要：
（1）训练数据的准备（校准集 + 微调数据），
（2）约为原始训练 5--10\% 的训练步数
（对于 70B 模型，这仍然意味着数千 GPU-hours），
（3）量化策略的超参数搜索（位宽、块大小、温度等）。
因此 QAT 通常只在以下场景中使用：
目标位宽为 INT4 或更低、
部署规模大到训练成本可以被摊薄、
或者精度要求特别严格（如医疗和金融应用）。
对于大多数部署场景，
PTQ（GPTQ 或 AWQ）的精度已经完全可以接受。

\subsection{量化的质量--大小权衡}

\begin{keybox}[量化位宽 vs.\ 精度损失经验法则]
\begin{itemize}[noitemsep]
  \item \textbf{FP8 / INT8}：近乎无损。
        大多数基准测试上精度损失 $<0.5\%$。
        数据中心部署的标准选择。
  \item \textbf{INT4（好的量化方法）}：1--3\% 的精度损失。
        对大多数实际应用可以接受。
        消费级 GPU 的最佳平衡点。
  \item \textbf{INT4（简单 RTN）}：3--8\% 的精度损失。
        质量明显下降，不推荐。
  \item \textbf{NVFP4（Blackwell 原生）}：1--2\% 精度损失，
        吞吐量达到 FP8 的 $2\times$。
        数据中心新一代标准。
  \item \textbf{INT3 / INT2}：显著退化。
        目前技术下仅适用于极端资源受限场景。
        研究活跃但尚未成熟。
\end{itemize}
\end{keybox}

\begin{keybox}[量化方法速查（2026 年初）]
\begin{itemize}[noitemsep]
  \item \textbf{数据中心 GPU（H100/H200）}：FP8，近零精度损失
  \item \textbf{Blackwell GPU（B200）}：NVFP4，$2.3\times$ 吞吐提升
  \item \textbf{消费级 GPU + 高吞吐}：AWQ + Marlin kernel
  \item \textbf{消费级 GPU + 兼容性}：GPTQ + BitBLAS kernel
  \item \textbf{CPU / Apple Silicon}：GGUF Q4\_K\_M（llama.cpp / Ollama）
  \item \textbf{微调}：QLoRA（4-bit 量化权重 + LoRA 适配器）
\end{itemize}
\end{keybox}

一个反直觉的发现：\textbf{量化不仅节省内存，还能提升速度}。
在 H200 上对 Qwen2.5-32B 的实测中，
AWQ 4-bit + Marlin kernel 的吞吐量达到 741 tok/s，
反而超过了 FP16 的 461 tok/s。
这是因为 LLM 推理是\textbf{内存带宽瓶颈}的，
瓶颈不是计算，而是从显存中读取权重。
更小的权重意味着更少的内存读取，因此更快。

\begin{table}[H]
\centering\small
\renewcommand{\arraystretch}{1.3}
\caption{INT4 量化方法在 H200 上的实测对比（Qwen2.5-32B-Instruct）}
\begin{tabularx}{\linewidth}{l c c c X}
\toprule
\rowcolor{figLightGray}
\textbf{方法} & \textbf{Perplexity} & \textbf{HumanEval} & \textbf{吞吐 (tok/s)} & \textbf{特点} \\
\midrule
FP16 基线 & 6.28 & 85.4\% & 461 & 参考基准 \\
GPTQ-INT4 & 6.42 & 82.9\% & 624 & Hessian 优化，通用性好 \\
AWQ-INT4 & 6.33 & 84.1\% & 741 (Marlin) & 激活感知，质量最佳 \\
GGUF Q4\_K\_M & 6.38 & 83.5\% & --- & CPU/Apple 首选 \\
BitsAndBytes & 6.51 & 81.7\% & 398 & 简单易用，QLoRA 标配 \\
\bottomrule
\end{tabularx}
\end{table}

\subsection{量化方法的选择决策}

面对多种量化方法，选择的核心考量是\textbf{目标硬件}和\textbf{精度要求}：

\textbf{GPTQ vs AWQ}：
两者都是 PTQ 方法，都针对 INT4 优化。
GPTQ 通过 Hessian 信息补偿量化误差，
对所有权重一视同仁地优化；
AWQ 则识别出关键权重给予优先保护。
在大多数基准上 AWQ 精度略优于 GPTQ，
且量化速度更快（AWQ 不需要计算完整的 Hessian 逆）。
但 GPTQ 的生态更成熟，在更多推理框架中有优化的 kernel 支持。

\textbf{GPU 量化 vs GGUF}：
GPTQ 和 AWQ 产生的模型需要 GPU 推理框架（vLLM、TensorRT-LLM）；
GGUF 格式面向 llama.cpp 生态，
支持 CPU、Apple Silicon Metal、CUDA 和 Vulkan 多种后端。
如果目标是在 MacBook 上运行模型，GGUF 是唯一实用的选择。
如果目标是数据中心高吞吐部署，AWQ + Marlin 是最优选。

\textbf{量化 Kernel 的重要性}不亚于量化算法本身。
一个被低估的事实是：
量化后的模型如果没有对应的高效 kernel，
反而可能\textbf{比 FP16 更慢}。
这是因为量化推理需要额外的解量化步骤，
将 INT4/INT8 权重转换回浮点数再做矩阵乘法。
如果解量化操作不够快，其开销会抵消甚至超过
减少内存读取带来的收益。

高效量化 kernel 的设计本质上是一个
\textbf{访存模式优化}问题。
Marlin~\cite{frantar2024marlin}（GPTQ/AWQ 的高性能 kernel）
通过三个关键技术实现了接近理论带宽上限的性能：
（1）\textbf{异步全局到共享内存拷贝}，
在 GPU 的 Tensor Core 做计算的同时，
异步加载下一批量化权重到 SRAM；
（2）\textbf{寄存器级解量化}，
INT4 到 FP16 的转换在寄存器中完成，
不需要写回共享内存再读取；
（3）\textbf{数据布局重排}，
量化权重的存储顺序经过精心设计，
使得 128-bit 的内存访问可以一次获取 32 个 INT4 值，
完美匹配 Tensor Core 的输入宽度。
这些底层优化使得 AWQ + Marlin 在 H200 上
达到了 FP16 理论带宽的 90\% 以上利用率，
量化不再有“速度税”。

BitBLAS 是另一个值得关注的量化 kernel 框架。
与 Marlin 针对特定量化方案手写 kernel 不同，
BitBLAS 使用\textbf{编译器自动生成}的方式
为任意位宽组合生成优化 kernel。
这种方法牺牲了极致性能（通常比手写 kernel 慢 10--20\%），
但大幅提升了灵活性，
支持 GPTQ、AWQ、QuIP\#~\cite{chee2024quip} 等多种量化方案，
以及 INT2、INT3、NF4 等非主流位宽。
对于需要快速试验不同量化配置的研究者，
BitBLAS 的“一键生成 kernel”能力极具价值。

\begin{whybox}[为什么量化后的模型速度有时反而更慢？]
一个常见的困惑：我把模型量化到 INT4 了，
为什么推理反而比 FP16 更慢？

原因通常有三个：
（1）使用了不匹配的 kernel，
例如在 NVIDIA GPU 上使用了为 CPU 优化的解量化路径；
（2）batch size 太大，
当 batch size 超过临界值（约 150--300），
推理从内存带宽瓶颈变为计算瓶颈，
量化减少内存读取的收益消失，
但解量化的额外计算成本依然存在；
（3）分组量化（group size）太小，
过小的分组大小（如 32）带来了过多的缩放因子，
解量化时需要频繁读取和应用不同的缩放因子，
反而增加了内存访问模式的不规则性。

经验法则是：
小 batch（$B < 64$）时 INT4 几乎总是更快，
中 batch（$64 < B < 256$）时取决于 kernel 质量，
大 batch（$B > 256$）时 FP8 通常是更好的选择
（因为 FP8 可以直接利用 Tensor Core 的原生 FP8 运算路径，
不需要解量化）。
\end{whybox}


\section{GGUF K-Quant：层级量化的艺术}

对于在消费级硬件上运行 LLM 的用户，
llama.cpp 的 GGUF 格式是事实标准。
其中最受欢迎的 Q4\_K\_M 格式使用了精巧的\textbf{层级量化}。

\subsection{为什么需要分块量化？}

最简单的量化方式是\textbf{per-tensor 量化}，
整个权重矩阵共享一个缩放因子。
问题是：如果矩阵中有个别极大的值，
缩放因子被迫设得很大，
导致大多数小值被量化到零或极少的几个级别，精度严重损失。

\textbf{Block quantization}（分块量化）
将权重矩阵切分为固定大小的块（block），
每个块独立计算缩放因子和零点。
块大小通常为 32--256 个权重。
这样每个块内的值分布更加均匀，
量化精度远好于全局量化。
代价是额外存储每个块的元数据（缩放因子和零点），
但这个开销通常只有 0.5--1 bit/weight。

\subsection{K-Quant 的两级结构}

GGUF 的 K-quant 采用两级量化结构：
\begin{enumerate}[noitemsep]
  \item 层的权重被切分为 \textbf{super-block}（如 256 个权重一组），
        每组配一个 FP16 缩放因子。
  \item 每个 super-block 内部包含多个 \textbf{sub-block}（如 32 个权重一组），
        各自独立计算缩放因子并量化权重。
  \item sub-block 的缩放因子本身被\textbf{二次量化}（6-bit），
        由 super-block 级别的 FP16 缩放因子覆盖。
\end{enumerate}

这种层级结构让量化误差被控制在更小的局部范围内，
避免了全局量化中「大值吃掉小值」的问题。
二次量化是一个巧妙的设计：
sub-block 的缩放因子本身也是数值，也可以被量化。
用 6-bit 量化缩放因子几乎不损失精度，
但节省了大量存储，
如果每个缩放因子都用 FP16 存储，
256 个权重中有 8 个 FP16 缩放因子（$8 \times 16 / 256 = 0.5$ bit/weight），
而二次量化将这个开销降低到约 0.25 bit/weight。

\subsection{K-Quant 变体}

GGUF 提供了一系列 K-Quant 变体，
命名规则是 \texttt{Q\{bits\}\_K\_\{size\}}：
\begin{itemize}[noitemsep]
  \item \texttt{Q2\_K}：2-bit 量化。模型极小但质量较差。
  \item \texttt{Q3\_K\_S / Q3\_K\_M / Q3\_K\_L}：3-bit 量化的小/中/大变体。
        后缀 S/M/L 指的不是块大小，而是\textbf{混合精度策略的激进程度}，
        L 变体对更多敏感层使用更高位宽。
  \item \texttt{Q4\_K\_S / Q4\_K\_M}：4-bit 量化。
        Q4\_K\_M 是最流行的选择，在质量和大小间取得了最佳平衡。
  \item \texttt{Q5\_K\_S / Q5\_K\_M}：5-bit 量化。接近 FP16 质量。
  \item \texttt{Q6\_K}：6-bit 量化。几乎无损。
  \item \texttt{Q8\_0}：8-bit 量化。完全无损，但大小是 Q4\_K\_M 的两倍。
\end{itemize}

\subsection{混合精度策略}

Q4\_K\_M 中的 “M” 代表 Medium 精度变体，
使用\textbf{混合精度策略}：
对质量影响最大的层（attention 的 Q/K/V 和 FFN 的 down\_proj）
使用 Q6\_K（6-bit），
其余层使用 Q4\_K（4-bit）。
这仅多出 $\sim$0.2 bits/weight 的存储成本，
但将 perplexity 增量从 0.1149 降低到 0.0535，
质量提升约 2$\times$。

这个设计背后的洞察是：
\textbf{不同层对量化误差的敏感度差异巨大}。
注意力层的 Q/K/V 直接影响「模型关注什么」，
量化误差会在 softmax 中被放大；
而 FFN 的 up\_proj 和 gate\_proj 对微小的数值误差
有更强的容忍度。
通过对敏感层分配更多比特，
可以用极小的额外成本获得显著的质量提升。

此外，模型的\textbf{第一层和最后几层}通常也更敏感：
第一层负责将 token embedding 投影到模型的内部表示空间，
最后几层负责将内部表示转换为 token 概率分布。
这两端的量化误差会直接影响输入理解和输出质量，
没有后续层来「纠正」误差。

\subsection{llama.cpp 与 GGUF 格式}

llama.cpp 之所以成为本地推理的事实标准，
有几个关键原因：
\begin{itemize}[noitemsep]
  \item \textbf{纯 C/C++ 实现}，无 Python 依赖，
        可以在任何平台（Linux、macOS、Windows、甚至手机）上编译运行。
  \item \textbf{GGUF 文件格式}是自包含的：
        一个文件包含模型权重、分词器、超参数和元数据。
        不需要额外的配置文件或依赖。
  \item \textbf{CPU 推理优化}：利用 AVX2/AVX-512（x86）
        和 NEON（ARM）指令集加速量化矩阵乘法。
  \item \textbf{GPU offloading}：可以将部分层放在 GPU 上，
        实现 CPU--GPU 混合推理，
        即使 GPU 显存不够放下整个模型，
        也能利用 GPU 加速最耗时的层。
  \item \textbf{Apple Silicon 原生支持}：
        通过 Metal API 利用 Apple M 系列芯片的 GPU 和 ANE（Neural Engine），
        在 MacBook 上实现高效推理。
\end{itemize}

Ollama 在 llama.cpp 之上提供了更友好的用户体验，
一条 \texttt{ollama run llama3} 命令就能下载和运行模型，
对非技术用户极为便利。
但 Ollama 在高并发场景下的性能远不及专业推理框架，
在 128 并发请求的压力测试中，
Ollama 的吞吐量仅为 484 tok/s，
约为 vLLM 或 SGLang 的 $1/10$。


\section{Flash Attention 与内存高效注意力}

标准的注意力计算需要将完整的 $n \times n$ 注意力矩阵
写入 GPU 的高带宽显存（HBM），
内存复杂度为 $O(n^2)$。
当序列长度 $n$ 达到数万甚至数十万时，
注意力矩阵本身就可能耗尽所有显存。

Flash Attention~\cite{dao2022flashattention}
通过\textbf{分块计算}（tiling）和\textbf{重计算}（recomputation）策略，
避免将完整的注意力矩阵写入 HBM，
将注意力计算的内存复杂度从 $O(n^2)$ 降低到 $O(n)$。
核心思想是：将 $Q$、$K$、$V$ 矩阵分成小块，
每次只在 GPU 的片上 SRAM（速度比 HBM 快约 $10\times$）中
计算一个子块的注意力，
逐块累积结果，始终不将完整的 $n \times n$ 矩阵写回 HBM。

Flash Attention 2~\cite{dao2023flashattention2}
进一步优化了 GPU 线程块的工作分配，
减少了非矩阵乘法运算（如 softmax 的归一化），
在 A100 上实现了接近理论峰值 72\% 的计算利用率。
Flash Attention 3~\cite{shah2024flashattention3}
针对 H100 的异步执行和 FP8 支持做了进一步优化。

这一技术同时加速了训练和推理，
是现代 LLM 部署的标配优化：
几乎所有主流推理框架（vLLM、SGLang、TensorRT-LLM）
和训练框架（Megatron-LM、DeepSpeed）都已集成 Flash Attention。

\subsection{Flash Attention 的工作原理详解}

要理解 Flash Attention 为什么有效，
需要深入 GPU 的\textbf{内存层次结构}。
现代 GPU 有三个主要的存储层级：
\begin{itemize}[noitemsep]
  \item \textbf{HBM}（高带宽显存）：
        容量大（H100: 80\,GB），带宽高（3.35\,TB/s），
        但延迟相对较高（$\sim$400 ns）。
        所有模型权重和 KV cache 存储在这里。
  \item \textbf{L2 Cache}：
        容量中等（H100: 50\,MB），带宽极高（$\sim$12\,TB/s），
        自动缓存 HBM 的热点数据。
  \item \textbf{SRAM}（片上共享内存）：
        容量极小（H100: 每个 SM 228\,KB，总计 $\sim$26\,MB），
        但带宽最高（$\sim$19\,TB/s），延迟最低（$\sim$30 ns）。
        程序员可以\textbf{显式管理}的唯一存储层。
\end{itemize}

标准注意力的问题在于它的内存访问模式：
计算 $\text{Softmax}(QK^T/\sqrt{d})$ 需要将整个 $n \times n$ 的注意力矩阵
\textbf{写入} HBM，然后再从 HBM 中\textbf{读回}来做 softmax 归一化，
最后再次写回归一化后的结果。
这个“读-写-读-写”的模式使得 HBM 成为瓶颈，
注意力计算本身的 FLOPS 很少（$O(n^2 d)$），
但 HBM 的 I/O 量是 $O(n^2)$。

Flash Attention 的核心技巧是\textbf{在线 softmax}（online softmax）。
传统 softmax 需要两遍遍历：
第一遍找到最大值（数值稳定性），
第二遍计算归一化的概率。
这要求整个注意力行都在内存中。
Milakov \& Gimelshein (2018) 提出了一种
\textbf{单遍}计算 softmax 的方法，
在扫描每个元素时\textbf{同时维护}
当前的最大值 $m$ 和累积指数和 $\ell$，
当遇到更大的值时回溯更新之前的累积量。

有了在线 softmax，Flash Attention 的算法变得清晰：
\begin{enumerate}[noitemsep]
  \item 将 $Q$ 矩阵分成行块（每块 $B_r$ 行），
        $K, V$ 分成列块（每块 $B_c$ 列）。
  \item 对 $Q$ 的每个行块，依次遍历所有 $K, V$ 列块：
        在 SRAM 中计算当前块对的 $QK^T$ 子矩阵，
        用在线 softmax 更新注意力权重，
        累积 $\text{softmax}(\cdot) \cdot V$ 的部分结果。
  \item 所有列块遍历完后，最终结果直接写回 HBM。
\end{enumerate}
整个过程中，$n \times n$ 的注意力矩阵
\textbf{从未在 HBM 中完整存在}，
它被逐块地在 SRAM 中临时计算然后丢弃。
HBM 的 I/O 量从 $O(n^2)$ 降到了 $O(n^2 d / M)$，
其中 $M$ 是 SRAM 大小。
对于典型的参数设置（$d = 128$, $M \sim 200$\,KB），
这意味着实际 I/O 减少了 $\sim$10$\times$。

\subsection{Flash Attention 的演进}

从 v1 到 v3，Flash Attention 的每一代
都针对当代 GPU 架构做了定向优化：

\textbf{Flash Attention 2} 的关键改进有两个。
第一，改变了\textbf{并行化维度}，
FA1 在 batch 和 head 维度并行，
FA2 在序列长度维度也并行，
减少了非矩阵乘法运算（如 softmax）的占比。
第二，优化了 warp 级别的工作分配，
让 warp 之间共享 $Q$ 块而不是 $K/V$ 块，
减少了 warp 间的同步需求。
这两项改进使得 FA2 在 A100 上达到了
理论峰值的 72\% 计算利用率（FA1 约 52\%）。

\textbf{Flash Attention 3}~\cite{shah2024flashattention3}
针对 H100 的三个硬件特性做了深度优化：
（1）\textbf{异步 warpgroup 操作}，
H100 的新 warpgroup-level 指令允许
一组 warp 在 Tensor Core 做矩阵乘法的同时，
另一组 warp 在 SRAM 中做 softmax 计算，
实现了计算和非计算操作的\textbf{重叠}；
（2）\textbf{FP8 注意力}，
利用 H100 的原生 FP8 Tensor Core，
在精度允许的场景下将注意力计算的
吞吐量提升 $2\times$；
（3）\textbf{更精细的 SRAM 双缓冲}，
在加载下一个 $K/V$ 块的同时
计算当前块的注意力，
完全隐藏 SRAM 加载延迟。
FA3 在 H100 上达到了约 740\,TFLOPS 的实际 BF16 注意力吞吐，
是理论峰值的 75\%：
在一个“非纯矩阵乘法”的计算中
达到这个利用率，体现了极致的硬件感知优化。

值得注意的是，Flash Attention 对\textbf{推理}的加速
主要体现在\textbf{prefill 阶段}和\textbf{长序列场景}。
在 decode 阶段（每次只有一个 query token），
注意力计算量极小（$O(n \cdot d)$，线性于序列长度），
Flash Attention 的优化效果有限。
Decode 阶段的瓶颈是\textbf{模型权重的读取}，
而非注意力计算，这也是为什么推测解码
和量化在 decode 阶段的加速效果更显著。


\section{KV Cache：空间换时间}

\subsection{为什么需要 KV Cache？}

LLM 的文本生成是自回归的，每次生成一个 token。
生成第 $t$ 个 token 时，
注意力机制需要计算它与前面所有 $t-1$ 个 token 的关系。
如果每次都重新计算所有位置的 $K$ 和 $V$，
生成一个长度为 $T$ 的序列需要
$O(T^2)$ 的注意力计算，因为第 $t$ 步需要 $O(t)$ 的计算。

\textbf{KV cache} 将之前位置计算过的 $K$ 和 $V$ 向量缓存起来，
每次只需计算新 token 的 $Q$，
与缓存的 $K, V$ 做注意力运算。
这将每步的新增计算量从 $O(t)$ 降到 $O(1)$（不含缓存读取），
总复杂度从 $O(T^2)$ 降到 $O(T)$。

\subsection{KV Cache 的内存开销}

KV cache 的大小计算：
\begin{equation}
  \mathrm{KV\ cache\ size} = 2 \times n_{\mathrm{layers}}
  \times n_{\mathrm{kv\_heads}} \times d_{\mathrm{head}}
  \times \mathrm{seq\_len} \times \mathrm{batch\_size}
  \times \mathrm{bytes}
\end{equation}

以 LLaMA 3-70B 为例（80 层、8 个 KV 头、128 维头维度），
在 BF16 精度下，\textbf{单个请求}处理 4096 个 token 的 KV cache：
\begin{equation}
  2 \times 80 \times 8 \times 128 \times 4096 \times 2\,\text{bytes}
  = 1.34\,\text{GB}
\end{equation}

如果上下文长度增加到 128K：
$2 \times 80 \times 8 \times 128 \times 131072 \times 2 \approx 42\,\text{GB}$，
单个请求的 KV cache 就超过了一张 A100 40GB 的全部显存。

当并发服务数百个请求时，KV cache 可能占用大部分 GPU 显存。
这使得 KV cache 的管理成为推理系统的核心挑战之一。

\subsection{PagedAttention}

vLLM~\cite{kwon2023vllm} 引入了 \textbf{PagedAttention}，
借鉴操作系统的虚拟内存分页机制，
将 KV cache 分成固定大小的「页」（page），
按需分配和回收。

传统的 KV cache 管理为每个请求预分配一块\textbf{连续}内存，
大小等于最大可能的序列长度。
问题在于：
\begin{itemize}[noitemsep]
  \item \textbf{内部碎片}：大多数请求不会用满预分配的空间。
        如果预分配 2048 个 token 的空间，
        但请求只生成了 200 个 token，
        90\% 的内存被浪费了。
  \item \textbf{外部碎片}：请求的生命周期不同，
        内存中出现大量不连续的空闲区域，
        无法分配给新请求。
\end{itemize}

PagedAttention 的解决方案与操作系统的虚拟内存如出一辙：
\begin{enumerate}[noitemsep]
  \item 将 GPU 显存分成固定大小的\textbf{物理页}（block），
        每页存储固定数量 token 的 KV 值（如 16 个 token）。
  \item 每个请求维护一个\textbf{页表}（page table），
        将逻辑位置映射到物理页。
  \item KV cache \textbf{不需要连续存储}，
        同一个请求的 KV 可以分散在不同的物理页中。
  \item 按需分配：请求生成新 token 时，才分配新的物理页。
\end{enumerate}

这种设计将内存浪费从 60--80\% 降低到接近 0\%，
使得 vLLM 可以在相同的 GPU 上服务 2--4$\times$ 更多的并发请求。

PagedAttention 还引入了一个操作系统中经典的优化技术：
\textbf{Copy-on-Write}（写时复制，COW）。
在 parallel sampling（同一 prompt 生成多个候选回答）
和 beam search 等场景中，
多个输出序列共享相同的\textbf{输入前缀}。
如果为每个序列都复制一份前缀的 KV cache，
内存开销会乘以并行数 $N$。
COW 的做法是：所有序列\textbf{共享}前缀对应的物理页，
每个序列的页表指向同一组物理页。
只有当某个序列需要\textbf{修改}（即追加新的 KV cache）时，
才为该序列复制一份新的物理页。

这种设计在 Best-of-N 采样中效果尤为显著。
假设从同一 prompt 并行生成 $N=16$ 个回答，
prompt 长度 4096 tokens，平均输出长度 256 tokens。
不使用 COW 时，KV cache 内存 $= 16 \times (4096 + 256) = 69632$ 个 token 的 KV。
使用 COW 后，KV cache 内存 $= 4096 + 16 \times 256 = 8192$ 个 token，
减少了 $88\%$。
这使得在单张 GPU 上进行大规模 Best-of-N 采样成为可能。

\begin{whybox}[PagedAttention 的性能开销在哪里？]
PagedAttention 引入了\textbf{间接寻址}，
每次读取 KV cache 时需要先查页表确定物理地址，
再跳转到对应的物理页。
这比连续内存的顺序访问多了一次间接寻址开销。

在实践中，这个开销被两个因素抵消：
（1）每个物理页内的 KV 仍然是\textbf{连续存储}的
（页大小通常为 16 tokens，足够一次 cache line 读取），
间接寻址只发生在页的边界处；
（2）PagedAttention 的 CUDA kernel 经过精心优化，
页表查找与 KV 读取\textbf{流水线化}，
在计算当前页的注意力时预取下一页的地址。
vLLM 的基准测试表明，
PagedAttention 相比连续 KV cache 的注意力计算开销
不超过 5\%，而内存利用率的提升带来了远超这个代价的收益。
\end{whybox}

\subsection{KV Cache 压缩}

除了更好的内存管理，还可以从\textbf{减少 KV cache 本身的大小}入手：

\begin{figure}[ht]
\centering
\begin{tikzpicture}[
  block/.style={draw, minimum width=12mm, minimum height=8mm, font=\small},
  physblock/.style={block, fill=figLightGray},
  reqA/.style={block, fill=figBlue!20},
  reqB/.style={block, fill=figRed!20},
  shared/.style={block, fill=figGreen!20},
  label/.style={font=\small\bfseries},
  >=Latex
]
% Physical memory blocks
\node[label] at (-1.5, 3.2) {GPU 显存（物理页）};
\foreach \i/\c in {0/shared, 1/reqA, 2/reqB, 3/shared, 4/reqA, 5/reqB, 6/physblock, 7/physblock} {
  \node[\c] (p\i) at ({mod(\i,4)*1.5}, {3 - floor(\i/4)*1.2}) {P\i};
}

% Page table for Request A
\node[label, text=figBlue] at (-1.5, 0.2) {请求 A 页表};
\node[reqA] (a0) at (0, 0.2) {→P0};
\node[reqA] (a1) at (1.5, 0.2) {→P3};
\node[reqA] (a2) at (3.0, 0.2) {→P1};
\node[reqA] (a3) at (4.5, 0.2) {→P4};

% Page table for Request B
\node[label, text=figRed] at (-1.5, -0.8) {请求 B 页表};
\node[reqB] (b0) at (0, -0.8) {→P0};
\node[reqB] (b1) at (1.5, -0.8) {→P3};
\node[reqB] (b2) at (3.0, -0.8) {→P2};
\node[reqB] (b3) at (4.5, -0.8) {→P5};

% COW annotation
\draw[->, thick, figGreen] (a0.north) to[out=90,in=-90] (p0.south);
\draw[->, thick, figGreen] (b0.north) to[out=90,in=-90] (p0.south);
\node[font=\footnotesize, text=figGreen, align=center] at (6.2, 0.2)
  {共享前缀\\(Copy-on-Write)};

% Free blocks annotation (offset to avoid overlap with P7)
\node[font=\footnotesize, text=figGray] at (6.0, 2.4) {空闲};
\draw[->, thin, figGray] (5.8, 2.2) -- (p6.north east);
\draw[->, thin, figGray] (5.8, 2.2) -- (p7.north east);
\end{tikzpicture}
\caption{PagedAttention 的虚拟内存布局示意。
请求 A 和请求 B 共享前缀对应的物理页 P0 和 P3（COW），
各自的生成内容存储在独立的物理页中。
KV cache 不需要连续存储，空闲页可以按需分配。}
\label{fig:paged-attention}
\end{figure}

\textbf{GQA}（Grouped-Query Attention）~\cite{ainslie2023gqa}
在模型架构层面减小 KV cache。
标准 MHA 中每个 query head 有独立的 K/V head；
GQA 让多个 query head \textbf{共享}同一组 K/V head。
LLaMA 3-70B 使用 64 个 query head 和 8 个 KV head（8:1 共享），
KV cache 直接缩小到 MHA 的 $1/8$。

GQA 的成功引发了一个自然的问题：
如果 8:1 的共享有效，\textbf{能否走得更极端？}
MQA（Multi-Query Attention）是 GQA 的极端情况：
所有 query head 共享\textbf{同一个} KV head（64:1 共享），
KV cache 缩小到 MHA 的 $1/64$。
PaLM~\cite{chowdhery2022palm} 和 Falcon~\cite{penedo2023refinedweb}
使用了 MQA，确实获得了极小的 KV cache。
但 MQA 的精度损失在某些任务上不可忽视
（$\sim$1--2\% on average），
因为所有 head 共享同一组 KV 表示，
模型无法为不同的注意力模式
维护不同的“记忆视角”。

GQA 是 MHA 和 MQA 之间的\textbf{连续谱}：
\begin{itemize}[noitemsep]
  \item MHA：$n_{\text{kv}} = n_{\text{q}}$，精度最高，KV cache 最大。
  \item GQA-$g$：$n_{\text{kv}} = n_{\text{q}} / g$，
        精度与 MHA 接近（$g \leq 8$ 时几乎无损），
        KV cache 缩小 $g$ 倍。
  \item MQA：$n_{\text{kv}} = 1$，KV cache 最小，但精度有损。
\end{itemize}
2024--2026 年的主流模型几乎都选择了 GQA（$g = 4$ 或 $g = 8$），
因为它在 KV cache 大小和精度之间提供了最佳平衡。
例外是 DeepSeek，它选择了更激进的 MLA 方案，
用更复杂的架构换取更极致的 KV cache 压缩。

\textbf{MLA}（Multi-head Latent Attention）
是 DeepSeek-V2/V3~\cite{deepseekv3} 提出的更激进的方案。
MLA 将 K 和 V 联合压缩到一个低秩表示 $c_t$：
\begin{equation}
  c_t = W^{DKV} [k_t; v_t], \quad
  [k_t; v_t] = W^{UKV} c_t
\end{equation}
其中 $c_t$ 的维度远小于原始 K/V 的拼接。
只需要缓存 $c_t$ 而不是完整的 K 和 V，
KV cache 大小可以压缩到原来的 $1/10$ 甚至更小。

\textbf{KV cache 量化}是另一条路径：
将 KV cache 从 BF16 量化到 INT8 甚至 INT4。
最新的研究表明，KV cache 对量化的耐受度
甚至高于权重，INT4 KV cache 的精度损失通常可以忽略不计。
这是因为 KV cache 中的数值分布比权重更集中、更平滑。
在 vLLM 中，开启 FP8 KV cache 量化只需一个配置参数，
即可将 KV cache 内存减半，几乎不影响输出质量。

KV cache 量化与权重量化的一个微妙差别在于\textbf{在线性质}。
权重量化可以在部署前离线完成（PTQ），
有充分的时间搜索最优的缩放因子；
KV cache 量化必须\textbf{在线}完成，
每个新 token 的 KV 计算出来后立即量化存储。
这意味着 KV cache 量化的缩放因子
只能基于当前已见到的 KV 值范围计算，
无法利用未来 token 的信息。

实践中，这种在线量化通常比离线量化\textbf{更容易}。
原因是 KV 值的分布在同一层内相当稳定，
前 100 个 token 的 KV 值范围
与前 1000 个 token 的范围相差不大。
这使得基于前几个 token 估计的缩放因子
对后续 token 同样适用。
TensorRT-LLM 的 FP8 KV cache 实现
使用了\textbf{per-channel 动态缩放}，
每个 head 维度独立计算缩放因子，
在每步解码时用指数移动平均更新。
这种策略在 INT4 KV cache 下也保持了
不到 0.5\% 的精度损失~\cite{hooper2024fp8}。

\textbf{KV cache 驱逐}（eviction）策略
在 KV cache 超出显存容量时发挥作用。
核心思想是：并非所有历史 token 的 KV 都同等重要。
通过注意力分数来评估每个 token 的「重要性」，
驱逐最不重要的 token 的 KV cache。
典型的策略包括：
\begin{itemize}[noitemsep]
  \item \textbf{Window attention}：只保留最近 $W$ 个 token 的 KV，
        丢弃更早的 token。简单但可能丢失关键的早期信息。
  \item \textbf{H2O}（Heavy-Hitter Oracle）：
        保留累积注意力分数最高的 token（“heavy hitters”）
        加上最近的窗口。这些 heavy hitters 只占少数 token
        但集中了大部分的注意力权重。
  \item \textbf{EvicPress}~\cite{feng2024evicpress}：
        联合优化压缩和驱逐决策，
        在多层存储（GPU HBM、CPU DRAM、SSD）之间
        动态调度 KV cache，最小化平均生成延迟。
\end{itemize}

\textbf{跨层 KV cache 共享}（Cross-layer KV sharing）
是一个新兴的研究方向。
观察表明，相邻 transformer 层的 KV 表示
具有很高的余弦相似度，
这意味着可以让多层共享同一组 KV cache，
而不是每层独立存储。
CommonKV~\cite{wang2025commonkv} 利用 SVD 分解
实现相邻层参数共享，
在不需要重新训练的情况下将 KV cache 压缩 2--4$\times$。
这种方法与量化和驱逐策略\textbf{正交}，
可以叠加使用，实现更极致的压缩。

\subsection{Prefix Caching}

在生产环境中，许多请求共享相同的\textbf{系统 prompt}
（system prompt），例如「你是一个有帮助的 AI 助手……」。
如果每个请求都重新计算系统 prompt 的 KV cache，
是巨大的浪费。

\textbf{Prefix caching} 将共享前缀的 KV cache 缓存起来复用：
第一个请求计算完系统 prompt 的 KV cache 后，
后续所有使用相同系统 prompt 的请求直接复用，
只需要计算用户特定输入部分的 KV cache。

对于典型的 API 服务（系统 prompt 占输入的 50--80\%），
prefix caching 可以将 prefill 计算量减少 50--80\%，
TTFT 显著降低。
vLLM 支持自动 prefix caching，
SGLang 则通过 RadixAttention 实现了更灵活的\textbf{树形}前缀复用，
不仅支持线性前缀，还能处理分支对话树结构。

\begin{whybox}[Prefix caching 为什么对多轮对话如此重要？]
在多轮对话中，每一轮都需要将完整的对话历史作为输入。
第 $N$ 轮对话的输入包含前 $N-1$ 轮的所有内容。
不使用 prefix caching 时，每一轮都需要重新计算
所有历史内容的 KV cache，
第 10 轮对话的 prefill 成本是第 1 轮的 10 倍。

使用 prefix caching 后，每一轮只需计算\textbf{新增}内容的 KV，
之前所有轮次的 KV cache 直接复用。
在 SGLang 的基准测试中，
多轮对话场景下 prefix caching 将 TTFT 降低了 5--10$\times$。
\end{whybox}

Prefix caching 的实现有几种不同的粒度：
\begin{itemize}[noitemsep]
  \item \textbf{Exact prefix matching}（精确前缀匹配）：
        只有当两个请求的 token 序列\textbf{完全相同}（逐 token 匹配）时
        才命中缓存。这是最简单也最安全的实现：
        任何一个 token 的差异都不会导致错误复用。
        vLLM 的默认 prefix caching 使用这种方式。
  \item \textbf{Hash-based matching}（哈希匹配）：
        将前缀按固定大小的块（如 16 tokens）分割，
        对每个块计算哈希值。
        两个请求只要共享前 $k$ 个块的哈希相同，
        就可以复用前 $k$ 个块的 KV cache。
        SGLang 的 RadixAttention 就是基于哈希的：
        基数树的每个边标签是一个 token 块的哈希。
  \item \textbf{Semantic matching}（语义匹配，实验性）：
        使用嵌入模型判断两个前缀是否“语义等价”，
        即使 token 序列不同但含义相同时也命中缓存。
        这在理论上能大幅提升命中率
        （例如同义改写的系统 prompt），
        但带来了\textbf{正确性风险}，
        语义相似的前缀可能导致不同的 KV cache，
        错误复用会产生无声的输出质量下降。
        目前仅在研究阶段。
\end{itemize}

\textbf{Prompt caching} 是 prefix caching 的商业化版本。
Anthropic 和 OpenAI 的 API 都支持 prompt caching：
对于重复使用的长前缀（如 RAG 的文档上下文），
第一次请求会缓存该前缀的 KV cache，
后续请求的 API 价格按缓存命中计费
（通常是标准价格的 $1/10$）。
这使得“长上下文 + 多次查询”的使用模式
在经济上变得极具吸引力，
例如上传一份 100 页的合同后多次提问，
只有第一次需要全价的 prefill 计算，
后续所有问题都按缓存价格计费。


\section{推测解码：小模型辅助大模型}

\subsection{核心洞察}

推测解码~\cite{leviathan2023speculative}
利用了一个观察：大部分 token 是「容易」的：
大模型和小模型对它们的预测一致。
只有少数「难」的 token 需要大模型的判断力。

Decode 阶段的瓶颈不是计算而是内存带宽：
不管生成一个 token 还是同时验证 $k$ 个 token，
读取模型权重的成本几乎相同
（因为 $k$ 通常很小，额外的计算量可以忽略）。
推测解码正是利用了这个特性。

\subsection{算法流程}

\textbf{流程}：
\begin{enumerate}[noitemsep]
  \item 用一个小的草稿模型（draft model，如 1B 参数）
        快速自回归生成 $k$ 个候选 token。
        由于模型很小，这一步很快。
  \item 将这 $k$ 个 token 一次性送入大模型做\textbf{并行验证}。
        大模型在一次前向传播中同时计算所有 $k$ 个位置的概率分布。
  \item 从前到后逐个检查：
        如果草稿模型的 token 被接受，继续检查下一个；
        如果某个位置被拒绝，
        从大模型在该位置的分布中重新采样，
        丢弃后续所有草稿 token。
\end{enumerate}

关键在于第 2 步：验证 $k$ 个 token 只需要一次前向传播
（因为可以并行计算），而不是 $k$ 次。
如果草稿模型的接受率是 80\%，
每次验证 5 个 token，
平均每次前向传播产生 4 个有效 token，
接近 $4\times$ 的加速。

\subsection{数学保证}

推测解码最重要的性质是：
它是\textbf{数学上无损的}，
最终输出的概率分布与直接使用大模型\textbf{完全等价}。

接受/拒绝机制使用了一种修正的\textbf{拒绝采样}（rejection sampling）。
设草稿模型在位置 $i$ 的分布为 $q(x)$，
目标大模型的分布为 $p(x)$。
对于草稿 token $x_i$：

\textbf{接受概率}：
\begin{equation}
  P(\text{accept } x_i) = \min\!\left(1,\, \frac{p(x_i)}{q(x_i)}\right)
\end{equation}

\textbf{拒绝后的重采样分布}：
当 $x_i$ 被拒绝时，从修正分布中重新采样：
\begin{equation}
  p'(x) = \frac{\max(0,\, p(x) - q(x))}{\sum_{x'} \max(0,\, p(x') - q(x'))}
\end{equation}

\begin{whybox}[为什么这个过程保证输出分布不变？]
对于任意 token $x$，被最终选中的总概率等于：
\begin{align}
  P(\text{select } x) &= q(x) \cdot \min\!\left(1, \frac{p(x)}{q(x)}\right) + P(\text{reject}) \cdot p'(x) \\
  &= \min(q(x), p(x)) + \left(1 - \sum_{x'} \min(q(x'), p(x'))\right) \cdot \frac{\max(0, p(x) - q(x))}{Z}
\end{align}
其中 $Z = \sum_{x'} \max(0, p(x') - q(x')) = 1 - \sum_{x'} \min(q(x'), p(x'))$。
代入化简后恰好等于 $p(x)$。

这个证明的关键在于：
$\sum \min(q, p) + \sum \max(0, p-q) = \sum p = 1$，
被接受的概率质量和被拒绝后重采样的概率质量
加起来恰好等于目标分布 $p$。
\end{whybox}

这意味着推测解码是\textbf{免费的加速}：
不牺牲任何输出质量。

\subsection{树形推测：从线性到指数级候选}

标准推测解码是\textbf{线性}的：
草稿模型生成一条长度为 $k$ 的 token 序列，
大模型逐位验证。
如果第 $i$ 个位置被拒绝，后面所有 token 都被丢弃。
这意味着一次拒绝就浪费了 $k - i$ 个 token 的生成成本。

\textbf{树形推测}（Tree-based speculation）
将候选从一条链扩展为一棵\textbf{树}：
在每个位置，草稿模型不是只生成一个 token，
而是生成多个候选 token，形成分支。
这棵推测树可以被大模型在\textbf{一次前向传播}中验证，
关键洞察是注意力计算天然支持
不同序列位置之间的并行计算，
只需要为每个候选 token 构造正确的\textbf{注意力掩码}
（确保每个候选只能看到它在树中的祖先节点）。

考虑一个具体例子：
在深度 3、分支因子 2 的推测树中，
总共有 $2 + 4 + 8 = 14$ 个候选 token。
线性推测同样的深度只有 3 个候选。
14 个候选中至少有一条完整路径被接受的概率
远高于单条链被完全接受的概率。

SpecInfer~\cite{miao2024specinfer}
将树形推测与\textbf{多草稿模型}结合，
使用多个不同的草稿模型各生成一棵推测子树，
合并为一棵更大的推测树。
不同草稿模型的“多样性”使得推测树覆盖更广的 token 空间，
提高了至少一条路径被接受的概率。
在 7B 目标模型上，SpecInfer 实现了 1.5--2.8$\times$ 的加速，
超过单草稿线性推测的 1.3--2$\times$。

EAGLE-2 的树形推测策略尤为精巧。
它不是使用固定的树结构，
而是根据上下文的\textbf{可预测性}动态调整树的形状：
\begin{itemize}[noitemsep]
  \item 高确定性位置（草稿模型置信度 $>$ 0.9）：
        只生成 1 个候选，不分支，因为几乎肯定会被接受。
  \item 中等确定性位置（置信度 0.5--0.9）：
        生成 2--3 个候选，轻度分支。
  \item 低确定性位置（置信度 $<$ 0.5）：
        生成 4--8 个候选，重度分支，
        因为需要更多候选来提高命中概率。
\end{itemize}
这种\textbf{自适应树结构}使得每次验证的 token 数量
可以根据文本的局部特性灵活变化。
在代码生成（结构化、高可预测性）中，树几乎退化为长链；
在创意写作（开放性、低可预测性）中，树变得宽而浅。

\begin{whybox}[树形推测为什么不需要更多的验证计算？]
直觉上，14 个候选 token 应该比 3 个需要更多验证计算。
但实际上，树形推测的验证成本只比线性推测高约 10--30\%，
而不是 $14/3 \approx 4.7\times$。

原因是 LLM decode 阶段是\textbf{内存带宽瓶颈}的。
验证 3 个 token 和验证 14 个 token 都需要
读取完整的模型权重（140\,GB for 70B），
差异只在于矩阵乘法的“宽度”，
从 3 个 query 变为 14 个 query。
由于 GPU 的计算能力远超内存带宽，
额外的 11 个 query 的计算几乎“免费”地被
闲置的计算单元消化了。
只有当候选 token 数量超过临界 batch size 时
（$\sim$150--300），验证成本才会显著增加。
在实践中，推测树的总节点数通常控制在 20--60 个，
远在计算瓶颈之下。
\end{whybox}

\subsection{加速比分析}

推测解码的加速比取决于两个关键因素：
\textbf{接受率} $\alpha$ 和\textbf{草稿长度} $k$。

设每次草稿生成 $k$ 个 token，每个 token 的接受率为 $\alpha$。
每次验证循环中，期望被接受的 token 数为：
\begin{equation}
  \mathbb{E}[\text{accepted tokens}] = \sum_{i=1}^{k} \alpha^i = \frac{\alpha(1 - \alpha^k)}{1 - \alpha}
\end{equation}
加上拒绝后重采样的 1 个 token，
每次循环产生的有效 token 数约为 $\frac{1 - \alpha^{k+1}}{1 - \alpha}$。

设草稿模型的推理时间为 $t_d$，大模型的推理时间为 $t_v$，
加速比为：
\begin{equation}
  \text{Speedup} = \frac{\frac{1 - \alpha^{k+1}}{1 - \alpha}}{k \cdot t_d + t_v} \cdot t_v
\end{equation}

当 $t_d \ll t_v$（草稿模型远小于目标模型）且 $\alpha$ 较高时，
加速比接近 $\frac{1}{1-\alpha}$。
例如 $\alpha = 0.8$ 时理论加速比约 $5\times$，
实际中由于草稿模型开销和内存管理，
通常实现 2--3$\times$ 的端到端加速。

\subsection{草稿模型选择与进化}

草稿模型的选择直接影响接受率和加速比：

\textbf{同族小模型}：最常用。
例如 Llama 3 8B 为 Llama 3 70B 做草稿。
词表和训练分布相似，接受率通常 70--85\%。
NVIDIA 在 H200 上展示了
用同族小模型实现 3.6$\times$ 吞吐量提升的结果。

\textbf{EAGLE 系列}~\cite{li2024eagle}：
将目标模型最后一层的隐状态
输入一个轻量级的自回归头来生成草稿。
EAGLE 考虑了 token 间的依赖关系，
生成的草稿通常比独立预测头更准确。

EAGLE-2 引入了\textbf{动态草稿长度}，
根据上下文的可预测性自适应调整
每次验证的候选 token 数量。
在「容易预测」的上下文中生成更多候选（如 8--10 个），
在「难预测」的上下文中减少候选（如 2--3 个），
避免浪费验证计算。

EAGLE-3~\cite{li2025eagle3} 进一步突破了性能瓶颈。
它放弃了 EAGLE 的特征预测范式，
改为直接预测 token，
并通过\textbf{多层特征融合}（training-time test 技术）
从目标模型的多个中间层提取信息。
这使得草稿模型能充分受益于训练数据规模的增加，
EAGLE 和 EAGLE-2 在增加训练数据时提升有限，
EAGLE-3 则表现出持续的改善。
在推理模型（如 DeepSeek-R1、QwQ）上，
EAGLE-3 实现了 2--4$\times$ 的加速。

\textbf{Medusa}~\cite{cai2024medusa}：在目标模型的最后一层之上
添加多个轻量级\textbf{预测头}（prediction head），
每个头预测不同位置的下一个 token。
多个头并行工作，类似于同时生成多个候选。
优点是不需要单独的草稿模型，
但需要额外训练预测头，
且由于各头独立预测，无法捕捉 token 间的依赖。

\textbf{Self-speculative decoding}（自推测解码）：
使用目标模型自身的\textbf{前几层}作为草稿模型。
不需要额外加载一个小模型，节省显存。

LayerSkip~\cite{elhoushi2024layerskip} 是这一方向的代表。
通过在训练中使用特殊的 dropout 策略
（早期层的 dropout 率低，后期层的 dropout 率高），
使得模型在只使用前几层时就能产生合理的预测。
推理时，使用前 $L/3$ 层作为草稿，
完整的 $L$ 层作为验证。
由于草稿和验证共享同一模型的权重，
无需额外的模型内存，
在内存受限场景下特别有吸引力。

CLaSp~\cite{chen2025clasp}（In-Context Layer Skip）
进一步改进了自推测解码，
不需要专门的训练，
通过分析上下文来动态决定跳过哪些层，
使得任何现有模型都可以直接使用自推测解码。

\textbf{HeteroSpec} 代表了推测解码的最新趋势，\textbf{自适应策略}。
通过输出分布的\textbf{熵}来判断当前上下文的难度：
低熵（高确定性）上下文使用更激进的推测（更多候选 token），
高熵（高不确定性）上下文使用保守的推测（更少候选）。
使用决策树将推测路径分为「快车道」和「慢车道」，
实现了 4.24$\times$ 的加速，超过 EAGLE-3。

\begin{keybox}[推测解码方法速查（2026 年初）]
\begin{itemize}[noitemsep]
  \item \textbf{同族小模型}：最成熟，生态最好。
        vLLM 和 TensorRT-LLM 原生支持。2--3$\times$ 加速。
  \item \textbf{EAGLE-3}：质量最高的轻量级草稿。
        需要训练草稿头。2--4$\times$ 加速。
  \item \textbf{Medusa}：无需独立草稿模型。
        多头并行预测。1.5--2.5$\times$ 加速。
  \item \textbf{LayerSkip / CLaSp}：自推测，零额外内存。
        适合内存受限场景。1.5--2$\times$ 加速。
  \item \textbf{HeteroSpec}：自适应策略，当前最快。
        需要熵估计开销。3--4$\times$ 加速。
\end{itemize}
\end{keybox}

实际加速通常在 2--3$\times$ 之间。
加速比取决于接受率，
接受率越高，每次验证产生的有效 token 越多。
对于代码生成等「可预测性高」的任务，加速比更大；
对于创意写作等「随机性高」的任务，加速比较小。
在推理模型（reasoning models）的 thinking token 生成中，
推测解码的加速效果尤为显著，
因为推理过程中的许多步骤是结构化的、可预测的。

\subsection{推测解码的适用性分析}

推测解码并非在所有场景下都有效。
理解它\textbf{何时有效}、\textbf{何时无效}
对于部署决策至关重要。

\textbf{推测解码效果好的场景}：
\begin{itemize}[noitemsep]
  \item \textbf{代码生成}：接受率通常 80--90\%。
        代码的语法结构高度可预测，
        关键词、缩进、括号匹配等模式
        小模型和大模型的预测高度一致。
        实测加速比 2.5--4$\times$。
  \item \textbf{推理模型的 thinking tokens}：接受率 75--85\%。
        推理过程中大量出现“Let me think...”
        “So the answer is...”等模板化结构，
        以及重复的数学公式格式。
        DeepSeek-R1 配合 EAGLE-3 可达 2--3$\times$ 加速。
  \item \textbf{长文档摘要}：接受率 70--80\%。
        摘要生成大量复用原文中的词汇和短语。
  \item \textbf{翻译}：接受率 65--80\%。
        特别是在术语密集的技术翻译中，
        专业术语的翻译在不同大小的模型间高度一致。
\end{itemize}

\textbf{推测解码效果差的场景}：
\begin{itemize}[noitemsep]
  \item \textbf{创意写作}：接受率低至 40--55\%。
        故事创作、诗歌等任务中，
        词汇选择的“正确答案”不唯一，
        大模型和小模型的偏好差异大。
        低接受率使得加速比降到 1.2--1.5$\times$，
        在考虑草稿模型的开销后可能接近无加速。
  \item \textbf{高温度采样}（$T > 1.0$）：
        高温度增加了输出的随机性，
        降低了大小模型预测的一致性。
        推测解码的数学保证仍然成立，
        但接受率大幅下降。
  \item \textbf{小 batch size + 计算密集型模型}：
        当 decode 已经\textbf{不是}内存带宽瓶颈时
        （例如大 batch 或长序列的 attention 计算），
        推测解码的基本假设（验证 $k$ 个 token $\approx$ 验证 1 个 token）
        不再成立，草稿模型的开销无法被隐藏。
  \item \textbf{多模态生成}（图像 token、音频 token）：
        非文本模态的 token 分布与文本差异巨大，
        同族小模型的草稿质量通常很差。
\end{itemize}

一个值得强调的\textbf{实践陷阱}是
推测解码对\textbf{batch size}的敏感性。
在单请求（$B = 1$）场景下，
推测解码的加速比最大，
因为 decode 完全是内存带宽瓶颈的，
额外的验证计算几乎“免费”。
但随着 batch size 增大，两个因素削弱了推测解码的收益：
（1）验证 $k$ 个 token 的额外计算不再可忽略，
batch 越大，每步的计算量越接近 GPU 上限，
额外的候选 token 开始竞争计算资源；
（2）草稿模型也需要处理大 batch，其开销不再微不足道。
经验法则是：$B < 16$ 时推测解码几乎总是有益的，
$16 < B < 64$ 时取决于具体配置，
$B > 64$ 时推测解码的收益通常被 continuous batching
的优化所取代，此时 GPU 利用率已经足够高，
推测解码的边际收益趋近于零。

\subsection{Speculative Speculative Decoding（SSD）}

2026 年 3 月，Stanford、Princeton 与 Together AI 的联合团队
（Tri Dao 参与共同作者）提出了
\textbf{Speculative Speculative Decoding（SSD）}，
将推测解码中的「推测」和「验证」两个阶段\textbf{并行化}。

传统推测解码是串行的：
草稿模型生成候选 $\to$ 目标模型验证 $\to$
接受/拒绝 $\to$ 草稿模型再次生成。
验证阶段目标模型工作时，草稿模型处于空闲状态。
SSD 的核心洞察是：
在目标模型进行验证的同时，
草稿模型可以\textbf{预测}验证可能的结果，
并为每种可能的结果\textbf{预先准备}下一轮的草稿 token。
当验证结果返回时，
直接选取对应的预生成草稿继续，无需等待。

这种「推测的推测」策略
在 H200 上实现了 \textbf{3.6$\times$} 的吞吐量提升，
显著超过传统推测解码的 2--3$\times$。

\textbf{QSpec} 是与 SSD 互补的另一项创新：
它为推测解码的草稿模型和验证模型
设计了\textbf{互补的量化方案}。
传统做法对两个模型使用相同的量化策略，
但 QSpec 发现联合激活-权重量化
在多步推理场景中会导致性能退化。
通过为草稿模型和验证模型分别优化量化配置，
QSpec 在不损失推测解码数学无损性的前提下
进一步提升了端到端速度。


\section{Continuous Batching：提升服务吞吐}

\subsection{Static Batching 的问题}

传统的推理服务（static batching）
将一批请求一起处理，等所有请求都生成完毕后才开始下一批。
问题是：不同请求的生成长度差异很大：
一个请求生成 10 个 token 就结束了，
另一个可能需要 2000 个 token。
短请求必须等待长请求完成，GPU 利用率很低。

假设一个 batch 有 8 个请求，
最短的生成 50 个 token，最长的生成 2000 个 token。
7 个短请求在前 50--200 步就完成了，
但它们的「位置」一直被占据到第 2000 步。
GPU 在 batch 后期有 80--90\% 的计算能力处于空闲状态。

\subsection{Iteration-Level Scheduling}

\textbf{Continuous batching}（又称 iteration-level batching）
解决了这个问题。
核心思想是：在\textbf{每一步}（每生成一个 token）
都重新审视当前 batch 的组成，
\begin{enumerate}[noitemsep]
  \item 检查是否有请求刚刚生成了 EOS token（结束标记）。
        如果有，将它从 batch 中移除。
  \item 检查等待队列中是否有新请求。
        如果有空位，立刻将新请求加入 batch。
  \item 对当前 batch 中的所有请求执行一步 decode。
\end{enumerate}

GPU 永远在满负荷工作，
不存在「等待最慢请求」的空闲时间。
在高并发场景下，
continuous batching 可以将吞吐量提升 2--5$\times$。

\begin{corebox}[Continuous Batching 的量化分析]
我们可以精确量化 static batching 的浪费。
设 batch 中有 $B$ 个请求，第 $i$ 个请求的生成长度为 $L_i$。
Static batching 的总计算量（以 decode 步数衡量）是 $B \times \max_i L_i$，
所有请求都等到最长的那个完成。
实际有效计算量只有 $\sum_i L_i$。
GPU 利用率为：
\begin{equation}
  \eta_{\text{static}} = \frac{\sum_i L_i}{B \times \max_i L_i}
\end{equation}
如果生成长度服从均匀分布 $U[1, L_{\max}]$，
期望利用率为 $\mathbb{E}[\eta] \approx \frac{1}{2} \cdot \frac{B+1}{B}$，
即使 batch 再大，也只有约 50\% 的利用率。
更糟糕的是，实际的生成长度分布通常是\textbf{重尾}的
（少数请求生成极长的输出），
这使得利用率进一步降低到 30--40\%。

Continuous batching 通过及时替换已完成的请求，
理论上可以将利用率推向 100\%
（受限于调度开销和 KV cache 管理的实际损耗，
通常达到 85--95\%）。
这 2--3$\times$ 的利用率提升\textbf{直接转化}为
相同硬件下的吞吐量提升，
这就是为什么 continuous batching
是推理框架最重要的基础优化之一。
\end{corebox}

\subsection{Chunked Prefill}

Continuous batching 还有一个微妙的问题：
当一个新请求加入 batch 时，它需要先经过 prefill 阶段
（处理其完整的 prompt）。
如果 prompt 很长（比如 32K tokens），
这次 prefill 会阻塞整个 batch 的 decode 步骤，
导致其他正在 decode 的请求出现\textbf{卡顿}（ITL 突增）。
在生产环境中，这种 ITL 突增会导致用户感知到明显的“打字卡顿”，
严重影响交互体验。
实测表明，一个 32K token 的 prefill 请求
可以使同一 batch 中其他请求的 ITL
从正常的 20--30\,ms 暴涨到 200--500\,ms，
超过人类感知阈值 (100\,ms) 数倍。

\textbf{Chunked prefill}~\cite{agrawal2024sarathi} 解决了这个问题：
将长 prompt 切成固定大小的块（chunk），
每步只处理一个块，
与同一步中的 decode token 一起组成一个\textbf{混合 batch}。
prefill 分散在多步中完成，
不会阻塞其他请求的 decode。

\textbf{最优 chunk 大小}的选择需要在三个因素间权衡：
\begin{itemize}[noitemsep]
  \item \textbf{ITL 稳定性}：chunk 越小，每步中 prefill 计算占比越低，
        decode 请求受到的干扰越小。极端情况下 chunk 大小为 1，
        则完全退化为逐 token prefill，ITL 完全平滑但 TTFT 大幅增加。
  \item \textbf{TTFT}：chunk 越大，prefill 完成越快，
        TTFT 越低。但过大的 chunk 又回到了阻塞 decode 的问题。
  \item \textbf{GPU 利用率}：chunk 大小应与 GPU 的计算能力匹配。
        在 H100 上，经验最优的 chunk 大小为 512--2048 tokens，
        使得每步的 prefill 计算量与 decode batch 的计算量相当，
        GPU 在两类计算间保持平衡。
\end{itemize}

Sarathi-Serve~\cite{agrawal2024sarathi} 对此进行了系统研究，
提出了\textbf{stall-free batching}的概念，
通过精心选择 chunk 大小，
保证每步的总计算时间波动不超过 10\%。
具体策略是将 chunk 大小设为
当前 decode batch 中 token 数量的函数：
decode batch 越大（GPU 计算负载越重），
prefill chunk 越小。
Sarathi-Serve 在 A100 集群上的实验表明，
stall-free batching 可以在\textbf{不损失吞吐量}的前提下
将 P99 ITL 降低 3--5$\times$。

Chunked prefill 与 continuous batching 的交互也值得注意。
在每一步的迭代中，调度器需要决定：
（1）为多少个新请求的 prefill chunk 分配计算资源，
（2）当前 decode batch 中有多少请求在等待。
这是一个在线调度问题：
过多的 prefill chunk 会挤占 decode 带宽（ITL 增加），
过少则浪费 GPU 的并行计算能力（TTFT 增加）。
vLLM 和 SGLang 都支持 chunked prefill，
SGLang 通过其 \textbf{Chunked Pipeline Parallelism}
将 chunked prefill 扩展到多 GPU 流水线并行场景，
在 1M token 长上下文的 prefill 中将 TTFT 降低了 81\%。

\subsection{Prefix-Aware Batching}

在支持 prefix caching 的系统中，
调度器可以进一步优化：
将共享相同前缀的请求\textbf{分组在同一个 batch} 中。
这样可以最大化 prefix cache 的命中率，
减少重复的 KV cache 计算。
朴素的 FCFS（先来先服务）调度完全忽略了请求间的前缀关系，
导致缓存频繁被不相关的请求驱逐，
即使大量请求共享同一个系统 prompt，
交错调度也会使缓存命中率低于 30\%。

SGLang~\cite{zheng2024sglang} 的 RadixAttention
维护了一棵\textbf{基数树}（radix tree）
来管理所有已缓存的前缀。
基数树的每个节点代表一段 token 序列的 KV cache，
从根到叶的路径对应一个完整的 prompt 前缀。
当新请求到来时，调度器在基数树中查找最长匹配前缀，
只需要计算未匹配的部分。
这种树形结构天然支持多种前缀共享模式：
\begin{itemize}[noitemsep]
  \item \textbf{线性前缀共享}：多个请求共享同一系统 prompt。
        这是最常见的场景，API 服务中所有请求通常有相同的系统指令。
  \item \textbf{分支前缀共享}：多轮对话中，
        不同分支的对话共享前 $N$ 轮的历史。
        基数树可以在分叉点上自然分支，避免重复缓存公共前缀。
  \item \textbf{Few-shot 前缀共享}：few-shot prompting 中，
        多个请求共享相同的示例但有不同的查询。
        基数树将共享示例的 KV cache 保留在父节点，
        每个查询只需计算自己独特的部分。
\end{itemize}

\textbf{缓存命中率}是 prefix-aware batching 的核心指标。
影响命中率的关键因素包括：
缓存容量（可用 GPU 显存）、请求的前缀分布
（是否集中在少数系统 prompt 上）、
以及\textbf{驱逐策略}。
常用的驱逐策略有三种：
LRU（Least Recently Used）最简单但可能驱逐高频前缀；
LFU（Least Frequently Used）保留热门前缀但对突发流量适应慢；
前缀长度加权的 LRU 优先保留较长的前缀
（因为长前缀的重计算成本更高）。
在典型的 API 服务场景中
（10--20 个不同的系统 prompt，每个服务数千请求），
prefix-aware scheduling 可以将缓存命中率从 30\% 提升到 85\% 以上，
TTFT 降低 2--5$\times$。

Prefix-aware batching 与 disaggregated serving 的交互
也是一个值得关注的设计维度。
在分离式架构中，prefill 节点计算 KV cache 后
需要将其传输到 decode 节点。
如果 prefill 节点维护了高效的前缀缓存，
大量请求的 KV cache 可以直接从缓存中获取
而不需要重新计算，
这不仅减少了 prefill 节点的计算负载，
也减少了 prefill-to-decode 的 KV cache 传输量
（因为 decode 节点也可以缓存常用前缀的 KV）。
vLLM 的 disaggregated serving 实现中，
prefix caching 与 KV cache 迁移是协同工作的：
如果某个前缀的 KV cache 已经在 decode 节点本地，
则完全跳过 prefill 阶段和跨节点传输，
延迟降到毫秒级。


\section{推理服务架构}

\subsection{推理框架对比}

\begin{table}[H]
\centering\small
\renewcommand{\arraystretch}{1.3}
\caption{主流推理框架对比（2026 年初）}
\begin{tabularx}{\linewidth}{l X}
\toprule
\rowcolor{figLightGray}
\textbf{框架} & \textbf{核心特性} \\
\midrule
\textbf{vLLM} & PagedAttention, Continuous Batching, Chunked Prefill, Prefix Caching,
  Speculative Decoding, 支持 disaggregated serving。Python-first，生态最丰富，
  社区最活跃（1969+ 贡献者）。
  适合数据中心通用部署。
  H100 吞吐约 12,500 tok/s（Llama 3.1 8B）。 \\
\textbf{SGLang} & RadixAttention（树形 prefix caching），Chunked Pipeline Parallelism
  （1M token TTFT 降低 81\%），结构化生成（grammar-constrained decoding）。
  H100 吞吐约 16,200 tok/s，比 vLLM 快 29\%。
  多轮对话场景再快 10--20\%。
  适合复杂 prompt 工程和 DeepSeek 模型部署。 \\
\textbf{TensorRT-LLM} & NVIDIA 原生优化, FP8/NVFP4 支持, In-flight Batching,
  最高单卡性能。C++ 核心，配置复杂但性能极致。
  适合 NVIDIA GPU 上的高性能生产部署。 \\
\textbf{LMDeploy} & TurboMind C++ 引擎，
  H100 吞吐约 16,100 tok/s，与 SGLang 持平。
  量化模型部署性能最优。
  适合需要极致量化推理的场景。 \\
\textbf{llama.cpp} & CPU/混合推理, GGUF 量化, Metal/CUDA/Vulkan 加速。
  纯 C/C++，无依赖，单文件模型格式。
  消费级硬件和边缘设备首选。 \\
\textbf{MLC-LLM} & 跨平台部署（iOS、Android、浏览器）。
  WebGPU 支持，可以在浏览器中运行 LLM。
  适合移动端和嵌入式场景。 \\
\bottomrule
\end{tabularx}
\end{table}

选择推理框架时需要考虑几个维度：
\begin{itemize}[noitemsep]
  \item \textbf{目标硬件}：NVIDIA GPU → vLLM / SGLang / TensorRT-LLM；
        Apple Silicon → llama.cpp；移动端 → MLC-LLM。
  \item \textbf{性能 vs.\ 易用性}：TensorRT-LLM 性能最高但配置复杂；
        vLLM 平衡了两者；SGLang 在多轮对话和 DeepSeek 上最优；
        llama.cpp 和 Ollama 最简单。
  \item \textbf{特殊需求}：结构化输出 → SGLang；
        浏览器部署 → MLC-LLM；多模态 → vLLM。
  \item \textbf{成本}：SGLang 比 vLLM 吞吐高 29\%，
        百万日请求场景下月省约 \$15,000 GPU 成本。
\end{itemize}

\begin{whybox}[vLLM vs SGLang：两种不同的设计哲学]
vLLM 和 SGLang 是 2025--2026 年数据中心部署中
最常被比较的两个框架。
它们的设计哲学有根本性的差异。

\textbf{vLLM 的核心是通用性}。
它的 PagedAttention 机制对所有模型架构透明，
OpenAI 兼容的 API 接口使得迁移成本极低。
vLLM 拥有最大的社区（1969+ 贡献者），
几乎所有新发布的开源模型都会在第一时间适配 vLLM。
它支持多模态模型推理（图文、视频）、
推测解码、disaggregated serving 等全套功能。
但 vLLM 的 Python 调度器在高并发下
可能成为瓶颈，
每一步的调度决策都需要 Python 解释器参与。

\textbf{SGLang 的核心是结构化推理}。
它的 RadixAttention 不仅仅是一种 prefix caching 方案，
更是一种\textbf{编程模型}，
开发者可以用 SGLang 的 DSL（Domain-Specific Language）
描述复杂的推理流程，
框架自动优化 KV cache 共享和调度。
SGLang 的调度器部分用 C++ 实现，
在高并发下开销更低。
在 DeepSeek 系列模型（V3/R1）上，
SGLang 的性能优势尤为明显，
因为 MoE + MLA 的组合需要精细的内存管理，
RadixAttention 的树形缓存在这方面更有优势。

实际选择建议：
如果部署的是\textbf{通用 API 服务}
（多种模型、多模态、需要丰富的社区支持），选 vLLM；
如果部署的是\textbf{特定的 DeepSeek/Qwen 模型}，
或需要\textbf{结构化生成}
（JSON schema 约束、grammar-constrained decoding），选 SGLang；
如果追求\textbf{极致单卡性能}且能接受更高的配置复杂度，
选 TensorRT-LLM。
\end{whybox}

\textbf{NVIDIA Triton Inference Server}
值得单独讨论，因为它不是一个推理\textbf{引擎}，
而是一个推理\textbf{编排平台}。
Triton 本身不包含 LLM 推理的核心优化
（如 PagedAttention、continuous batching），
而是作为后端编排层，
调度请求到底层的推理引擎（TensorRT-LLM、vLLM 等）。
Triton 的价值在于\textbf{生产级运维能力}：
多模型共享 GPU、模型版本管理、
自动扩缩容、A/B 测试、请求优先级队列等。
在大规模生产环境中，
Triton + TensorRT-LLM 的组合是 NVIDIA 推荐的标准方案。

\textbf{Hugging Face TGI}（Text Generation Inference）
定位于“零配置快速部署”。
对于 Hugging Face Hub 上的模型，
一条 \texttt{docker run} 命令即可启动推理服务。
TGI 内部使用了 Flash Attention 和 continuous batching，
性能在中等并发下与 vLLM 相当，
但在高并发（$>$100 并发）场景下吞吐量明显落后。
TGI 的主要优势在于与 Hugging Face 生态的深度集成，
模型下载、tokenizer 处理、Inference Endpoints 托管服务
都是开箱即用的。
对于快速验证和中小规模部署，TGI 是入门门槛最低的选择。

\subsection{Disaggregated Serving：分离 Prefill 和 Decode}

传统的推理框架将 prefill 和 decode 放在同一组 GPU 上执行。
但这两个阶段的计算特性截然不同：
\begin{itemize}[noitemsep]
  \item \textbf{Prefill} 是计算密集型，大量 token 并行计算，
        GPU 算力被充分利用。
  \item \textbf{Decode} 是内存带宽密集型，逐 token 生成，
        GPU 大部分时间在等待内存读取。
\end{itemize}

将两者混合在同一组 GPU 上会互相干扰：
长 prompt 的 prefill 会阻塞其他请求的 decode（ITL 突增），
大量 decode 请求又会延迟新 prompt 的 prefill（TTFT 增加）。

\textbf{Disaggregated serving}（分离式推理）
将 prefill 和 decode 分配到\textbf{不同的 GPU 节点}：
\begin{itemize}[noitemsep]
  \item \textbf{Prefill 节点}：处理输入 prompt，生成 KV cache，
        然后将 KV cache 传输到 decode 节点。
        适合使用计算能力强的 GPU。
  \item \textbf{Decode 节点}：接收 KV cache，逐 token 生成输出。
        适合使用内存带宽大的 GPU。
\end{itemize}

\begin{corebox}[DeepSeek 的分离式推理架构]
DeepSeek 在其 V3/R1 推理系统中
采用了大规模的分离式架构。
根据 2025 年 2 月 Open Source Week 披露的信息：

\begin{itemize}[noitemsep]
  \item Prefill 节点和 decode 节点的比例约为 1:4，
        decode 阶段占用了绝大部分的 GPU 资源。
  \item 跨节点的 KV cache 传输通过高速互联（NVLink/InfiniBand）完成，
        传输延迟被控制在毫秒级。
  \item 得益于 MLA 的 KV cache 压缩，
        需要传输的数据量远小于传统 MHA 架构，
        使得分离式架构的通信开销可控。
  \item 每个 H800 节点的吞吐量
        达到甚至超过同等活跃参数量（$\sim$37B）的稠密模型，
        这证明了 MoE + MLA + 分离式架构的组合效果。
\end{itemize}
\end{corebox}

Mooncake~\cite{qin2024mooncake} 进一步提出了
\textbf{KVCache-Centric}（以 KV cache 为中心）的分离式架构。
其核心创新是将 KV cache 视为独立的资源进行管理，
而不是绑定到特定的 GPU。
KV cache 可以存储在 GPU HBM、CPU DRAM 或 SSD 上，
调度器根据实时负载动态决定 KV cache 的存放位置。
Mooncake 还引入了\textbf{预测性早期拒绝}（Early Rejection）——
在请求到达时就预估系统是否能满足 SLO，
如果无法满足则提前拒绝，避免浪费计算资源。

vLLM v0.11 也引入了 disaggregated serving 支持，
配合 Wide Expert Parallel（Wide-EP）实现了
DeepSeek-R1 在 H200 集群上 2,200 tok/s/GPU 的吞吐。

\textbf{DistServe}~\cite{zhong2024distserve}
将分离式架构推向了更精细的粒度。
不同于 Mooncake 主要关注 KV cache 的调度，
DistServe 的核心贡献在于
\textbf{goodput-optimal}（有效吞吐量最优）的资源分配。
它将 SLO（Service Level Objective）
直接纳入调度器的目标函数——
不是简单地最大化吞吐量，
而是最大化满足延迟 SLO 的\textbf{有效请求数}。
这意味着当系统接近过载时，
DistServe 会主动降低 batch size 以保证延迟 SLO，
而不是盲目接受更多请求然后全部超时。

\textbf{Splitwise}~\cite{patel2024splitwise}
从\textbf{异构硬件}的角度重新审视分离式推理。
核心洞察是 prefill 和 decode 不仅计算特性不同，
它们的\textbf{最优硬件配置}也不同：
\begin{itemize}[noitemsep]
  \item Prefill 受益于更高的计算能力（FLOPS）——
        H100 SXM 的 989 TFLOPS 比 H100 PCIe 的 756 TFLOPS
        在 prefill 上快 30\%。
  \item Decode 受益于更高的内存带宽——
        H200 的 4.8\,TB/s HBM3e 带宽
        比 H100 SXM 的 3.35\,TB/s 快 43\%，
        decode 速度几乎线性加速。
\end{itemize}
Splitwise 的策略是将 prefill 集群配备“计算强”的 GPU
（如 H100 SXM），decode 集群配备“带宽强”的 GPU
（如 H200 或 HBM 版本的 A100），
让每种 GPU 只做它最擅长的工作。
这种异构部署在总 GPU 数量不变的前提下，
可以将总体有效吞吐提升 20--40\%。

分离式推理面临的关键挑战是
\textbf{KV cache 跨节点传输}的延迟。
Prefill 节点计算完 KV cache 后，
需要将其传输到 decode 节点才能开始生成。
对于一个 70B 模型、4096 token prompt 的 KV cache（约 1.34\,GB），
通过 InfiniBand（400\,Gbps $\approx$ 50\,GB/s）传输需要约 27\,ms。
这个传输延迟直接加到 TTFT 上。

缓解传输瓶颈的策略包括：
（1）\textbf{KV cache 压缩}——
MLA 将 KV cache 压缩到原来的 $1/10$，
传输量相应减少 $10\times$，延迟降到 $\sim$3\,ms；
（2）\textbf{KV cache 量化}——
将 KV cache 从 BF16 量化到 FP8 或 INT4，
传输量减半或减到 $1/4$；
（3）\textbf{流水线传输}——
不等 prefill 全部完成，
而是边计算边传输已完成层的 KV cache，
使得传输延迟与 prefill 计算\textbf{重叠}。
Mooncake 的实现结合了这三种策略，
将 KV cache 传输延迟控制在 5\,ms 以内——
对于大多数应用场景，这已经在用户感知阈值之下。

\subsection{多模态与 MoE 专用推理架构}

随着模型从纯文本走向多模态、从稠密走向 MoE，
推理架构也在相应进化。

\textbf{vLLM-Omni}（arXiv:2602.02204）
提出了\textbf{全解耦}的多模态推理架构，
支持\textbf{任意到任意}（any-to-any）的多模态模型服务——
文本、图像、视频、音频的输入输出可以任意组合。
其核心设计将不同模态的 prefill、decode 和生成阶段
完全解耦到独立的计算节点上，
每个模态的处理流水线可以独立扩缩容。
这解决了多模态模型在生产部署中的关键痛点：
不同模态的计算特征差异巨大
（视频 prefill 计算密集、文本 decode 带宽密集），
统一调度会导致严重的资源浪费。

\textbf{MoE-SpAc}（arXiv:2603.09983）
专门针对 MoE 模型在\textbf{异构边缘设备}上的推理优化。
核心思想是\textbf{推测性激活}（speculative activation）——
预测即将被路由到的专家，
提前将其权重从慢速存储加载到快速内存中。
这对于资源受限的边缘设备尤为重要：
MoE 模型的总参数量远大于设备内存容量，
必须按需加载专家，
而加载延迟是推理速度的主要瓶颈。
MoE-SpAc 通过路由预测机制将专家加载延迟隐藏在计算过程中，
使得大型 MoE 模型在边缘设备上的推理变得实用。

\subsection{MoE 模型推理的独特挑战}

MoE 模型的推理与稠密模型有本质性的差异，
这些差异使得传统推理优化不能直接套用。

\textbf{全专家加载问题}是最核心的挑战。
DeepSeek-V3 有 671B 总参数但只有 37B 活跃参数——
每次前向传播只激活 8/257 个专家。
但模型的\textbf{所有 671B 参数}都需要驻留在 GPU 显存中
（因为不同 token 路由到不同专家，
无法预知哪些专家会被使用）。
这意味着 MoE 模型的\textbf{内存需求}由总参数量决定（671B），
而\textbf{计算需求}由活跃参数量决定（37B）——
一个极端不平衡的“宽模型、浅计算”模式。

对于推理优化，这种不平衡有两个重要含义：

（1）\textbf{量化对 MoE 的收益更大}。
由于 MoE 模型的内存需求远超计算需求
（内存/计算比是稠密模型的 $\sim$18$\times$），
减少内存占用的任何手段都直接转化为
可以在更少的 GPU 上部署整个模型。
将 DeepSeek-V3 从 BF16（约 1.34\,TB）量化到 FP8（约 0.67\,TB）
或 INT4（约 0.34\,TB），
部署所需的 GPU 数量从 16$\times$ H100 降到 8$\times$ 或 4$\times$。
ExpertsInt8（前文提到的 AI21 方案）正是利用了这一点。

（2）\textbf{Expert parallelism}（专家并行）
是 MoE 推理的专用并行策略。
标准的张量并行（Tensor Parallelism）
将每一层的权重切分到多个 GPU 上，
每步都需要 all-reduce 通信。
Expert parallelism 则将不同的\textbf{专家}分配到不同 GPU，
只有被路由到远程 GPU 上的专家的 token
才需要跨 GPU 通信——
当路由是“本地”的（token 路由到同一 GPU 上的专家），
完全不需要通信。
DeepSeek 的 Wide-EP（Wide Expert Parallel）
将这一策略推向极致：
在 H200 集群上将 257 个专家分散到 64 个 GPU，
每个 GPU 负责约 4 个专家，
通过高速 NVLink 连接实现跨 GPU 的 token 调度。

（3）\textbf{Load balancing 对推理吞吐的影响}
在 MoE 推理中尤为关键。
如果某些专家被过度路由（“热门专家”），
持有这些专家的 GPU 成为瓶颈，
其他 GPU 空闲等待——
这与 static batching 中等待最慢请求的问题如出一辙。
DeepSeek-V3 在训练时使用了辅助负载均衡损失来缓解这一问题，
但推理时的 token 分布可能与训练时不同
（例如用户请求中数学 token 的比例可能偏高，
导致数学相关的专家被过度激活）。
动态负载均衡策略——
如将过热专家的部分 token 分流到“备用”专家——
是 MoE 推理系统需要解决的工程问题。

\subsection{边缘部署与端侧推理}

随着量化技术的成熟和硬件能力的提升，
在\textbf{边缘设备}上运行 LLM 正在从实验走向实用。
2026 年初，在 Apple M4 Pro MacBook 上运行
Llama 3.1 8B（Q4\_K\_M 量化，约 4.7\,GB）
可以达到 30--50 tok/s 的生成速度——
对于交互式聊天完全可用。
在 iPhone 16 Pro 上运行 Qwen2.5 3B（Q4\_0 量化）
也能达到 15--25 tok/s。

边缘推理与数据中心推理的优化重点截然不同：
\begin{itemize}[noitemsep]
  \item \textbf{功耗约束}：
        手机和笔记本的功耗预算（5--15\,W）
        远低于数据中心 GPU（300--700\,W）。
        优化目标从“最高吞吐”变为“每瓦特吞吐”。
        Apple M 系列芯片的统一内存架构
        在这方面有独特优势——
        CPU 和 GPU 共享同一内存池，
        消除了数据在 CPU DRAM 和 GPU HBM 之间的搬运开销。
  \item \textbf{内存容量}：
        消费级设备的内存（16--36\,GB 统一内存或显存）
        比数据中心 GPU（80\,GB HBM）小得多。
        模型必须量化到 4-bit 甚至 2-bit 才能装下。
        GGUF Q4\_K\_M 是当前的“甜蜜点”——
        一个 8B 模型仅占 4.7\,GB，
        在 16\,GB 内存的 MacBook Air 上舒适运行。
  \item \textbf{隐私}：
        边缘推理的一个重要优势是
        数据完全不离开用户设备——
        对于敏感场景（个人日记、医疗记录、法律文件），
        这比云端推理有本质性的隐私保障。
\end{itemize}

\textbf{NPU/ANE 推理}（Neural Processing Unit / Apple Neural Engine）
是边缘推理的新兴方向。
现代移动 SoC 几乎都集成了专用的 AI 加速单元：
Apple A17/M4 的 ANE（35 TOPS INT8）、
Qualcomm Hexagon（45 TOPS INT8）、
Google Tensor（高效 TPU 内核）。
这些加速单元在 INT8/INT4 推理上的
每瓦特效率远超 GPU——
但它们的编程模型受限，
不支持 Flash Attention 等复杂的内存优化。
如何在 NPU 上高效实现 LLM 推理
（特别是注意力计算和 KV cache 管理），
是边缘 AI 工程的前沿课题。

\subsection{Dual-Batch 与异步调度}

\textbf{Dual-batch}（双批次重叠）是 vLLM 和 SGLang 中的一项
前沿优化技术。
核心思想是将 prefill 和 decode 的计算
在时间维度上\textbf{重叠}——
在 GPU 执行一批 decode 的同时，
异步准备下一批 prefill 的数据。
这消除了 prefill 和 decode 之间的空闲间隙，
进一步提升了 GPU 利用率。

理解 dual-batch 的价值需要回到 GPU 的执行模型。
现代 GPU（H100/H200）支持多个 CUDA stream 并发执行，
但在标准推理流程中，
prefill 和 decode 被串行调度到同一个 stream 上——
prefill 完成后才开始 decode，decode 完成后才处理下一个 prefill。
这导致了两类空闲时间：
（1）prefill 结束到 decode 开始之间的\textbf{调度间隙}（$\sim$1--3\,ms）；
（2）一批 decode 完成后等待新请求 prefill 的\textbf{排空间隙}（$\sim$5--10\,ms）。
在高吞吐场景下，这些间隙累积起来可占总运行时间的 5--15\%。

Dual-batch 通过\textbf{双缓冲}（double-buffering）策略
消除这些间隙。
系统维护两个并发的 batch：
batch A 在 GPU 上执行 decode，
与此同时 batch B 的 prefill 数据在 CPU 上准备
（tokenization、KV cache 分配、页表构建）。
当 batch A 的 decode 完成时，
batch B 已经准备就绪可以立即开始 prefill，
batch A 的下一轮 decode 则在 batch B prefill 的同时准备。
GPU 和 CPU 形成了一个\textbf{流水线}，
两端永远不会同时空闲。
SGLang 的实现表明，dual-batch 在高并发场景下
可以将端到端吞吐量提升 10--20\%，
且不增加单请求延迟。

\textbf{异步调度}（async scheduling）
则将调度决策从 GPU 计算路径中移出。
传统调度器在每步计算前同步决定 batch 组成，
引入了几毫秒的调度延迟——
调度器需要遍历等待队列、检查 KV cache 可用性、
计算前缀匹配、决定 batch 组成，
这些操作都在 CPU 上执行但阻塞了 GPU 的计算。
异步调度器在 CPU 上持续运行，
在 GPU 计算当前步的同时准备好下一步的 batch，
使得调度开销几乎为零。

异步调度的实现需要处理一个微妙的\textbf{一致性问题}：
当调度器在 $t$ 时刻做出 $t+1$ 步的 batch 决策时，
$t$ 步的结果尚未产生——
有些请求可能在 $t$ 步就结束了（生成了 EOS），
但调度器在决策时不知道这一点。
解决方案是\textbf{乐观调度}：
假设当前 batch 中没有请求结束，
如果有请求实际结束了，
则在 $t+1$ 步开始前快速做一次轻量级的 batch 修正。
实践中，每步结束的请求数量通常很少（batch 的 1--5\%），
因此修正的开销可以忽略。
vLLM 自 v0.6 起默认开启异步调度，
报告称调度延迟从 $\sim$5\,ms 降低到 $<$0.1\,ms。


\section{模型蒸馏：大模型教小模型}

\subsection{知识蒸馏的基本原理}

知识蒸馏（Knowledge Distillation）的核心思想是：
用一个大的\textbf{教师模型}（teacher）的输出
来训练一个小的\textbf{学生模型}（student）。
相比让学生模型直接在原始数据上训练，
教师模型的输出包含了更丰富的信息——
不仅告诉学生「正确答案是什么」，
还通过概率分布中的\textbf{暗知识}（dark knowledge）
告诉学生「其他答案为什么不对，但有多接近正确」。

标准的知识蒸馏使用带温度的 softmax：
\begin{equation}
  p_i^{(\tau)} = \frac{\exp(z_i / \tau)}{\sum_j \exp(z_j / \tau)}
\end{equation}
其中 $\tau$ 是温度参数。$\tau > 1$ 时，
概率分布变得更「平滑」，
各类别之间的相对概率差异更加明显——
这正是暗知识的来源。

蒸馏损失函数通常是两部分的加权和：
\begin{equation}
  \mathcal{L} = \alpha \cdot \tau^2 \cdot \mathrm{KL}(p_{\text{teacher}}^{(\tau)} \| p_{\text{student}}^{(\tau)})
  + (1 - \alpha) \cdot \mathcal{L}_{\text{CE}}(y, p_{\text{student}})
\end{equation}
第一项让学生模仿教师的概率分布，
第二项保证学生在原始任务上的准确性。
$\tau^2$ 因子是为了补偿高温 softmax 的梯度缩放。

\subsection{LLM 蒸馏的实践}

在 LLM 时代，蒸馏有几种具体形式：

\textbf{Logits 蒸馏}：
让学生模型直接匹配教师模型在每个 token 位置的完整概率分布。
这是信息量最大的蒸馏方式，
但需要访问教师模型的 logits——
对于闭源模型（如 GPT-4）无法使用。

\textbf{Sequence-level 蒸馏}（或称 data distillation）：
用教师模型生成大量高质量的训练数据，
然后在这些数据上对学生模型做 SFT。
这是最常用的方式，因为只需要教师模型的\textbf{输出文本}，
不需要 logits。
Stanford Alpaca~\cite{taori2023alpaca} 首次展示了
用 GPT-3.5 生成的数据训练 LLaMA 7B 的效果——
学生模型在很多任务上达到了教师的 80--90\% 水平。

\textbf{Chain-of-Thought 蒸馏}：
教师模型不仅提供答案，还提供推理过程。
学生模型同时学习「怎么推理」和「推理出什么」。
这在数学和代码任务上效果尤为显著——
学生模型从教师的推理链中学到了\textbf{结构化的思考方式}。

\subsection{DeepSeek-R1 的蒸馏实验}

DeepSeek-R1~\cite{deepseekr1} 的蒸馏结果
是目前 LLM 蒸馏最令人印象深刻的案例。
DeepSeek 团队用 R1（671B MoE）作为教师模型，
生成了包含完整推理链的 800K 条训练数据，
然后对多个开源基础模型进行 SFT 蒸馏。

\begin{table}[H]
\centering\small
\renewcommand{\arraystretch}{1.3}
\caption{DeepSeek-R1 蒸馏模型性能}
\begin{tabularx}{\linewidth}{l c c c X}
\toprule
\rowcolor{figLightGray}
\textbf{蒸馏模型} & \textbf{AIME} & \textbf{MATH-500} & \textbf{LiveCodeBench} & \textbf{备注} \\
\midrule
R1-Distill-Qwen-1.5B & 28.9\% & 83.9\% & 16.9\% & 手机端可运行 \\
R1-Distill-Qwen-7B & 55.5\% & 92.8\% & 37.6\% & 消费级 GPU \\
R1-Distill-Qwen-14B & 69.7\% & 93.9\% & 53.1\% & 单卡 \\
R1-Distill-Qwen-32B & 72.6\% & 94.3\% & 57.2\% & 性能/成本最优 \\
R1-Distill-Llama-70B & 70.0\% & 94.5\% & 57.5\% & 接近教师水平 \\
\midrule
\textit{R1 教师（671B MoE）} & \textit{79.8\%} & \textit{97.3\%} & \textit{65.9\%} & \textit{参考} \\
\bottomrule
\end{tabularx}
\end{table}

几个关键发现：

\textbf{蒸馏比直接 RL 更高效}。
在相同基础模型上，
蒸馏得到的推理能力\textbf{显著优于}
直接用 GRPO 做 RL 训练。
例如 Qwen2.5-32B 直接 RL 训练在 AIME 上达到 47.0\%，
但用 R1 蒸馏后达到 72.6\%——差距巨大。
这表明强大教师模型的推理链
包含了通过 RL 从零学习难以获得的信息。

\textbf{小模型可以「学会」远超其自身能力的推理}。
R1-Distill-Qwen-7B 在 MATH-500 上达到 92.8\%，
而同规模的 Qwen2.5-7B（未蒸馏）只有约 70\%。
蒸馏将一个 7B 模型的数学推理能力
提升到了接近 70B 未蒸馏模型的水平。

\textbf{再蒸馏}（re-distillation）可以进一步提升质量。
Mobius Labs 的实验表明，
用更大的蒸馏模型（如 R1-Distill-70B）
对更小的模型（如 R1-Distill-14B）进行二次 logits 蒸馏，
在多个基准上还能获得 2--5\% 的提升，
且成本极低（仅 \$3--\$18 的 GPU 成本）。

\textbf{On-policy 蒸馏 vs Off-policy 蒸馏}
是蒸馏方法论中一个被低估的维度。
标准的 sequence-level 蒸馏是\textbf{off-policy}的——
学生模型在教师生成的数据上训练，
但这些数据是由\textbf{教师模型的策略}产生的，
不反映学生模型自身的行为模式。
这导致了一个分布不匹配问题：
训练时学生看到的是教师的推理风格，
推理时学生必须按自己的风格续写——
教师的“下一步”和学生自己会生成的“下一步”不同，
误差在长序列中累积。

\textbf{On-policy 蒸馏}让学生模型自己生成推理链，
然后用教师模型\textbf{评分和纠正}这些生成的链。
学生在自己的分布上学习，
教师只负责指出哪些是好的、哪些是坏的。
这类似于 DAgger（Dataset Aggregation）
在模仿学习中的作用——
让学习者在自己的策略产生的状态分布上获取专家反馈。
实验表明，on-policy 蒸馏在长序列推理任务上
比 off-policy 蒸馏提升 5--10\% 的准确率，
但代价是需要在训练中同时运行学生和教师模型，
计算成本是 off-policy 的 3--5$\times$。

\textbf{Progressive distillation}（渐进式蒸馏）
是另一种有前景的策略：
先蒸馏到一个“中间”大小的学生模型（如 32B），
然后用这个中间模型作为新的教师，
再蒸馏到更小的学生（如 7B）。
多级蒸馏链的每一步“跨度”更小，
每步的信息损失也更小。
这类似于将一个大跳转分解为多个小步——
70B → 32B → 14B → 7B 的三级蒸馏链
通常优于直接的 70B → 7B 一步蒸馏。
Mobius Labs 和 Microsoft 的多项实验
证实了渐进式蒸馏在 MATH 和 HumanEval 上
可以额外带来 3--8\% 的准确率提升。

\subsection{蒸馏的局限性}

蒸馏并非万能：
\begin{itemize}[noitemsep]
  \item \textbf{能力天花板}：学生模型的性能受限于教师模型。
        如果教师在某个领域不够好，学生也学不好。
  \item \textbf{遗忘问题}：过度蒸馏可能导致学生模型
        在教师未覆盖的领域退化。
  \item \textbf{分布差异}：如果学生模型的预训练数据分布
        与教师模型差异很大，蒸馏效果会打折扣。
  \item \textbf{推理能力 vs.\ 知识}：蒸馏擅长传递推理\textbf{能力}
        （如 CoT 推理的模式），但不擅长传递\textbf{知识}
        （如具体的事实记忆）——
        因为知识需要参数容量来存储，小模型的参数有限。
\end{itemize}

\subsection{蒸馏的工程实践}

蒸馏看似简单——生成数据然后训练——
但在工程实践中有许多容易被忽视的细节，
它们决定了蒸馏结果的最终质量。

\textbf{教师输出的采样策略}直接影响蒸馏数据的质量。
使用贪心解码（$T = 0$）生成的数据确定性最高、质量最好，
但\textbf{多样性极低}——对于同一个 prompt，
总是生成完全相同的输出。
学生模型会过拟合到教师的单一表达方式上。
使用较高温度（$T = 0.7$--$1.0$）增加多样性，
但引入了低质量的噪声样本。
实践中的最佳策略是\textbf{拒绝采样}（rejection sampling）——
对每个 prompt 用高温度生成多个候选（如 8--16 个），
用教师模型或独立的评估器筛选出最优的 1--2 个。
DeepSeek-R1 的蒸馏正是使用了这种策略——
从大量生成的推理链中只保留正确且格式规范的输出，
最终筛选出约 800K 条高质量训练数据。

\textbf{数据配比}是另一个关键工程决策。
纯蒸馏数据训练的学生模型
往往在教师未覆盖的领域严重退化——
例如用数学蒸馏数据训练的模型可能“忘记”日常对话能力。
解决方案是混合蒸馏数据和\textbf{通用 SFT 数据}：
蒸馏数据教会学生特定能力（推理、代码），
通用数据保持基础能力（对话、翻译、摘要）。
R1 的训练流程（阶段三）正是这种混合策略的体现——
800K 推理蒸馏数据 + 通用 SFT 数据的联合训练。

\textbf{学习率和训练步数}的选择影响蒸馏效率。
一个常见的错误是用过大的学习率或过多的训练步数——
这会导致学生模型\textbf{过拟合}教师的表面模式
（如特定的遣词造句），而非学到深层的推理能力。
经验法则是：蒸馏的学习率应为预训练的 $1/10$ 到 $1/100$，
训练步数不超过原始训练的 5\%。
Mobius Labs 的再蒸馏实验验证了这一点——
在 1000 步蒸馏和 10000 步蒸馏之间，
1000 步的泛化性能反而更好，
因为更少的步数避免了对教师特定表达的过拟合。

\textbf{蒸馏与 RL 的关系}是一个深层次的问题。
R1 的实验清晰地表明：
蒸馏出的推理能力在短期内远超从零 RL 训练。
但从长远看，蒸馏的能力\textbf{天花板低于 RL}——
因为蒸馏只能传递教师已有的能力，
而 RL 有可能发现教师未知的策略。
最有前景的路线可能是\textbf{蒸馏 + RL}的结合：
先用蒸馏快速启动（warm-start），
再用 RL 在蒸馏的基础上\textbf{超越}教师。
这种“站在巨人肩膀上再突破”的策略
在 AlphaGo 的演进中已经得到验证——
AlphaGo Master 先从人类棋谱中学习（蒸馏），
再通过自我对弈（RL）超越了人类最高水平。


\section{推理优化的整体思路}

回顾整个推理优化领域，
可以归纳出一个清晰的思维框架：
所有优化都在解决同一个根本问题——
\textbf{LLM 推理是内存带宽瓶颈的，而非计算瓶颈的}。

在自回归生成过程中，
每生成一个 token，需要读取\textbf{整个模型的权重}
（因为每个 token 都要经过所有层），
但只做少量的计算（一个 token 的矩阵乘法）。
GPU 的计算能力远未被充分利用——
大部分时间花在等待从显存读取权重。

理解了这个瓶颈，所有优化手段都变得清晰：
\begin{itemize}[noitemsep]
  \item \textbf{量化}减少每次需要读取的数据量
        （4-bit 权重是 16-bit 的 $1/4$ 大小），
        直接缓解带宽瓶颈。
  \item \textbf{Continuous batching} 增加每次读取权重后
        处理的 token 数量——
        读一次权重，同时处理 100 个请求的 token，
        提高计算与内存访问的比率。
  \item \textbf{推测解码}减少需要调用大模型的次数——
        小模型的权重更小，读取更快；
        大模型一次验证多个 token，减少总读取次数。
  \item \textbf{KV cache}避免重复读取和计算
        之前位置的 KV 值。
  \item \textbf{PagedAttention}减少 KV cache 的内存浪费，
        允许在固定显存下服务更多并发请求。
  \item \textbf{Prefix caching} 消除共享前缀的重复计算，
        在生产环境中大幅降低 TTFT。
  \item \textbf{Disaggregated serving} 分离 prefill 和 decode，
        让每种 GPU 只做它最擅长的工作。
  \item \textbf{蒸馏}从根本上减小模型大小，
        用一个 7B 模型替代 70B 模型来服务大部分请求。
\end{itemize}

从这个角度看，推理优化的未来方向也很明确：
任何能减少「每 token 需要读取的数据量」
或「增加每次读取后完成的有效工作」的技术，
都值得关注。
NVIDIA Blackwell 的 NVFP4 支持、
更大的片上 SRAM、更高带宽的 HBM4——
硬件和软件都在向同一个方向努力：
打破内存带宽的瓶颈。

\subsection{推理优化的经济学视角}

站在 2026 年初回顾，
推理优化领域的发展速度远超大多数人的预期。
以 GPT-4 等效能力的单位推理成本为基准，
2022 年 12 月（GPT-4 发布前夕）到 2026 年 3 月之间，
成本下降了约 $1000\times$——
从每百万 token 约 \$20 降至约 \$0.02。
这个下降速度甚至超过了摩尔定律
（两年 $4\times$ vs.\ 推理成本两年 $30\times$）。

成本下降的来源可以分解为几个独立的贡献：
\begin{itemize}[noitemsep]
  \item \textbf{硬件进步}（$\sim$3$\times$）：
        H100 → H200 → Blackwell 的 HBM 带宽和算力提升。
  \item \textbf{量化}（$\sim$4$\times$）：
        BF16 → FP8 → INT4 / NVFP4，减少内存占用和带宽需求。
  \item \textbf{Serving 优化}（$\sim$8$\times$）：
        Continuous batching、PagedAttention、chunked prefill、
        prefix caching 的累积效果。
  \item \textbf{推测解码}（$\sim$2--3$\times$）：
        EAGLE-3、HeteroSpec 等的 decode 加速。
  \item \textbf{模型蒸馏}（$\sim$10$\times$）：
        70B → 7B 蒸馏模型达到类似的推理质量。
  \item \textbf{MoE 架构}（$\sim$3$\times$）：
        671B 参数但只有 37B 活跃，
        有效参数利用率大幅提升。
\end{itemize}
这些因素是\textbf{可叠加}的——
量化 + serving 优化 + 推测解码 + 蒸馏
可以同时应用到同一个部署中，
成本下降的乘积效应解释了 $1000\times$ 的总降幅。

这个趋势对整个 AI 产业有深远影响。
当推理成本降到足够低时，
之前“因为太贵而不实用”的应用变得可行：
实时代码审查（每次提交都调用推理模型检查）、
全量客户对话分析（对每通客服录音做推理）、
持续性 Agent 工作流（24/7 运行的 AI 助手）。
推理成本的每一次 $10\times$ 下降
都会打开一个新的应用市场——
这就是为什么推理优化仍然是 LLM 领域
最具商业价值的技术方向之一。

\subsection{推理系统设计的权衡空间}

设计一个推理系统，
本质上是在一个\textbf{多维权衡空间}中做选择。
理解这些权衡，比掌握任何单一优化技术更重要。

\textbf{延迟 vs 吞吐量}是最基本的权衡。
小 batch size 给单个请求低延迟，
但 GPU 利用率低、吞吐量低、单位成本高。
大 batch size 提高吞吐量和 GPU 利用率，
但每个请求的 TTFT 和 ITL 都会增加。
生产系统需要根据 SLO 要求
找到这条权衡曲线上的最优点——
对于交互式聊天（SLO: TTFT $< 500$\,ms, ITL $< 50$\,ms），
batch size 需要控制在较小范围；
对于离线批处理（无 SLO），可以用最大 batch size 榨干 GPU。

\textbf{精度 vs 速度}是量化的权衡。
更低位宽意味着更快的推理和更低的成本，
但精度损失在某些任务上可能不可接受。
关键洞察是：精度损失不是全局均匀的——
同一个量化方案在翻译任务上几乎无损（$<0.1\%$），
在数学推理上可能有显著退化（$\sim$3\%）。
因此量化策略需要\textbf{按任务定制}，
而不是“一个位宽走天下”。

\textbf{内存 vs 计算}是推理系统的永恒主题。
KV cache 是典型的“空间换时间”——
用 GPU 显存存储历史 KV 来避免重复计算。
推测解码则是“计算换计算”——
用小模型的计算来减少大模型的调用次数。
Disaggregated serving 是“通信换优化”——
用跨节点通信来换取更好的 prefill/decode 隔离。
每种权衡的最优点都取决于具体的硬件配置、
模型架构和流量特性。

理解这些权衡后，
推理系统设计不再是一个“选择最好的技术”的问题，
而是一个“在约束条件下做最优组合”的工程优化问题。
最好的推理系统不是用了最多优化技术的系统，
而是针对其特定场景
在权衡空间中找到了\textbf{帕累托最优点}的系统。

\subsection{推理系统的可观测性}

在生产环境中部署推理系统后，
\textbf{持续监控}是保证服务质量的关键。
推理系统的监控维度远多于传统 Web 服务。

\textbf{性能指标}（必须实时监控）：
\begin{itemize}[noitemsep]
  \item \textbf{P50/P95/P99 TTFT 和 ITL}：
        中位数之外必须关注尾部延迟——
        P99 TTFT 可能是 P50 的 5--10$\times$
        （长 prompt 触发了非缓存 prefill）。
  \item \textbf{GPU 利用率}（计算和内存带宽分别监控）：
        如果计算利用率低但带宽利用率高，
        说明系统正确地运行在带宽瓶颈区间；
        如果两者都低，说明 batch size 不足。
  \item \textbf{KV cache 占用率}：
        当 KV cache 使用率接近 100\% 时，
        新请求必须等待或被拒绝——
        这通常先于 GPU 利用率饱和发生。
  \item \textbf{Prefix cache 命中率}：
        如果命中率从 80\% 突降到 30\%，
        可能意味着流量模式发生了变化
        （新的系统 prompt 被引入、
        缓存驱逐策略不适应当前负载）。
\end{itemize}

\textbf{质量指标}（定期采样评估）：
\begin{itemize}[noitemsep]
  \item \textbf{量化漂移}：
        量化模型在部署初期的精度测试可能正常，
        但在特定输入分布下出现退化——
        定期在代表性查询集上评估输出质量。
  \item \textbf{推测解码接受率}：
        如果接受率从 80\% 降到 60\%，
        可能意味着用户请求的类型发生了变化
        （从结构化的代码查询变为创意写作），
        此时应考虑动态关闭推测解码。
\end{itemize}

好的推理系统不是一次性部署完成的——
它需要持续的\textbf{观测 → 分析 → 优化}循环。
推理优化的下一个前沿不在单一技术的突破，
而在于\textbf{系统级的自适应能力}——
能够根据实时流量特性自动调整
量化精度、batch 策略、缓存策略和推测解码配置的系统，
才是真正的下一代推理基础设施。


%% ============================================================
%% NEW REFERENCES
%% ============================================================
% NEW REF: lin2024awq = Lin, Ji, et al. "AWQ: Activation-aware Weight Quantization for LLM Compression and Acceleration." MLSys 2024.
% NEW REF: xiao2023smoothquant = Xiao, Guangxuan, et al. "SmoothQuant: Accurate and Efficient Post-Training Quantization for Large Language Models." ICML 2023.
% NEW REF: li2024eagle = Li, Yuhui, et al. "EAGLE: Speculative Sampling Requires Rethinking Feature Uncertainty." ICML 2024.
% NEW REF: li2025eagle3 = Li, Yuhui, et al. "EAGLE-3: Scaling up Inference Acceleration of Large Language Models via Training-Time Test." arXiv:2503.01840, 2025.
% NEW REF: cai2024medusa = Cai, Tianle, et al. "Medusa: Simple LLM Inference Acceleration Framework with Multiple Decoding Heads." ICML 2024.
% NEW REF: elhoushi2024layerskip = Elhoushi, Mostafa, et al. "LayerSkip: Enabling Early Exit Inference and Self-Speculative Decoding." ACL 2024.
% NEW REF: chen2025clasp = Chen, Longze, et al. "CLaSp: In-Context Layer Skip for Self-Speculative Decoding." arXiv:2505.24196, 2025.
% NEW REF: feng2024evicpress = Feng, Shaoting, et al. "EvicPress: Joint KV-Cache Compression and Eviction for Efficient LLM Serving." arXiv:2512.14946, 2024.
% NEW REF: wang2025commonkv = Wang, Yixuan, et al. "CommonKV: Compressing KV Cache with Cross-layer Parameter Sharing." arXiv:2508.16134, 2025.
% NEW REF: qin2024mooncake = Qin, Ruoyu, et al. "Mooncake: A KVCache-Centric Disaggregated Architecture for LLM Serving." arXiv:2407.00079, 2024.
% NEW REF: agrawal2024sarathi = Agrawal, Amey, et al. "Taming Throughput-Latency Tradeoff in LLM Inference with Sarathi-Serve." OSDI 2024.
% NEW REF: zheng2024sglang = Zheng, Lianmin, et al. "SGLang: Efficient Execution of Structured Language Model Programs." arXiv:2312.07104, 2024.
% NEW REF: liu2024llmqat = Liu, Zechun, et al. "LLM-QAT: Data-Free Quantization Aware Training for Large Language Models." arXiv:2305.17888, 2024.
% NEW REF: kim2024squeezellm = Kim, Sehoon, et al. "SqueezeLLM: Dense-and-Sparse Quantization." ICML 2024. arXiv:2306.07629, 2024.
% NEW REF: frantar2024marlin = Frantar, Elias, et al. "Marlin: Mixed-Precision Auto-Regressive Parallel INference on GPUs." arXiv:2408.11743, 2024.
% NEW REF: chee2024quip = Chee, Jerry, et al. "QuIP#: Even Better LLM Quantization with Hadamard Incoherence and Lattice Codebooks." ICML 2024.
% NEW REF: miao2024specinfer = Miao, Xupeng, et al. "SpecInfer: Accelerating Generative Large Language Model Serving with Tree-based Speculative Inference and Verification." ASPLOS 2024.
% NEW REF: zhong2024distserve = Zhong, Yinmin, et al. "DistServe: Disaggregating Prefill and Decoding for Goodput-optimized Large Language Model Serving." OSDI 2024.
% NEW REF: patel2024splitwise = Patel, Pratyush, et al. "Splitwise: Efficient Generative LLM Inference Using Phase Splitting." ISCA 2024.
